\documentclass[12pt]{article}
\usepackage[utf8]{inputenc}
\usepackage[spanish]{babel}
\usepackage{amsmath, amssymb, amsthm}
\usepackage{geometry}
\usepackage{tikz}
\geometry{a4paper, margin=2.5cm}
\usepackage{hyperref}
\usepackage{xcolor}
\usepackage{amsmath}

\DeclareMathOperator{\cond}{cond}


\newtheorem{teorema}{Teorema}
\newtheorem{lema}{Lema}
\newtheorem{definicion}{Definición}
\newtheorem{proposicion}{Proposición}


\title{\textbf{Definiciones Formales}\\Elementos de Cálculo Numérico}

\begin{document}
\maketitle

\section*{Unidad 1 — Complejidad y Análisis Algorítmico}

\subsection*{1. Complejidad Computacional}
La \textbf{complejidad computacional} estudia los recursos necesarios
para la ejecución de un algoritmo en función del tamaño $n$ de su entrada.

Los principales recursos considerados son:

\begin{itemize}
    \item \textbf{Complejidad temporal:} cantidad de operaciones elementales
    necesarias para ejecutar el algoritmo.
    
    \item \textbf{Complejidad espacial:} cantidad de memoria necesaria para
    ejecutar el algoritmo.
\end{itemize}

Si $T(n)$ representa el costo temporal de un algoritmo para una entrada
de tamaño $n$, el análisis de complejidad estudia el comportamiento
asintótico de $T(n)$ cuando
\[
n \longrightarrow \infty.
\]


\subsection*{2. Cotas Asintóticas}
Sean
\[
f,g:\mathbb{N}\longrightarrow\mathbb{R}_{\geq 0}.
\]

Las \textbf{cotas asintóticas} permiten describir el orden de crecimiento
de $f(n)$ respecto de $g(n)$ cuando $n\to\infty$.
\begin{itemize}
    \item \textbf{Cota Superior $\mathcal{O}(g(n))$:} Representa el peor caso. \\
    $\mathcal{O}(g(n)) = \{ f(n) : \exists c > 0, n_0 > 0 \text{ tal que } 0 \le f(n) \le c \cdot g(n) \quad \forall n \ge n_0 \}$
    \item \textbf{Cota Inferior $\Omega(g(n))$:} Representa el mejor caso. \\
    $\Omega(g(n)) = \{ f(n) : \exists c > 0, n_0 > 0 \text{ tal que } 0 \le c \cdot g(n) \le f(n) \quad \forall n \ge n_0 \}$
    \item \textbf{Cota Ajustada $\Theta(g(n))$:} Comportamiento exacto. \\
    $\Theta(g(n)) = \mathcal{O}(g(n)) \cap \Omega(g(n))$
\end{itemize}

\subsection*{3. Árbol de Secuencia (Recursión)}
Diagrama estructural que representa el despliegue de una ecuación de recurrencia. Cada nodo detalla el costo iterativo de un subproblema en un nivel específico, permitiendo visualizar cómo se distribuye el trabajo de dividir y combinar la información hasta alcanzar los casos base (las hojas).

\subsection*{4. Teorema Maestro}
El \textbf{Teorema Maestro} permite determinar el comportamiento
asintótico de recurrencias de la forma
\[
\boxed{
T(n)=a\,T\!\left(\frac{n}{b}\right)+f(n)
}
\]
donde
\[
a\geq 1,\qquad b>1,
\]
y $f(n)$ representa el costo del trabajo realizado fuera de las
llamadas recursivas.

La recurrencia describe un problema que se divide en $a$ subproblemas,
cada uno de tamaño $n/b$.

El término fundamental de comparación es
\[
\boxed{
n^{\log_b a},
}
\]
que corresponde al orden de crecimiento de la cantidad total de hojas
del árbol de recurrencia.

El Teorema Maestro determina $T(n)$ comparando $f(n)$ con
$n^{\log_b a}$.

\subsection*{Casos del Teorema Maestro}

Sea
\[
T(n)=a\,T\!\left(\frac{n}{b}\right)+f(n),
\qquad a\geq 1,\quad b>1.
\]

\subsubsection*{Caso 1: $f(n)$ domina por debajo}

Si existe una constante $\varepsilon>0$ tal que
\[
\boxed{
f(n)=O\!\left(
n^{\log_b a-\varepsilon}
\right),
}
\]
entonces el costo de los niveles internos es asintóticamente menor
que el costo de las hojas, y
\[
\boxed{T(n)=\Theta\!\left(n^{\log_b a}\right).}
\]

\subsubsection*{Caso 2: equilibrio}

Si
\[
\boxed{
f(n)=\Theta\!\left(n^{\log_b a}\right),
}
\]
entonces $f(n)$ tiene el mismo orden de crecimiento que el costo total
de las hojas. En este caso,
\[
\boxed{
T(n)=\Theta\!\left(
n^{\log_b a}\log n
\right).
}
\]

\subsubsection*{Caso 3: $f(n)$ domina por encima}

Si existe una constante $\varepsilon>0$ tal que
\[\boxed{f(n)=\Omega\!\left(n^{\log_b a+\varepsilon}\right),}
\]
y además se cumple la \textbf{condición de regularidad}, es decir,
existe una constante $c<1$ tal que, para $n$ suficientemente grande,
\[
\boxed{
a\,f\!\left(\frac{n}{b}\right)\leq c\,f(n),
}
\]
entonces el costo de la raíz domina asintóticamente al de los niveles
inferiores y
\[
\boxed{
T(n)=\Theta(f(n)).
}
\]
\section*{Aritmética de Punto Flotante}

\subsection*{6. Números de Máquina y Sistema}
Un sistema de números de máquina es un subconjunto estrictamente finito y discreto de los números racionales que un hardware puede representar. Está determinado unívocamente por cuatro parámetros: $\beta$ (base), $m$ (límite de dígitos de la mantisa), $e_{min}$ y $e_{max}$ (rango del exponente).

\subsection*{7. Desarrollo en Base y Mantisa}
Todo número real $x$ se aproxima mediante un formato normalizado estricto:
\begin{equation}
    fl(x) = \pm (0.d_1 d_2 \dots d_m) \cdot \beta^e
\end{equation}
Donde la \textbf{mantisa} es el bloque $(0.d_1 d_2 \dots d_m)$ que almacena las cifras significativas. La normalización exige que el primer dígito fraccionario sea distinto de cero ($d_1 \neq 0$).

\subsection*{8. Propiedades Fundamentales}
\begin{itemize}
    \item \textbf{Finitud:} Al tener restricciones en $m$ y $e$, la cantidad total de números representables es discreta y limitada.
    \item \textbf{Densidad no uniforme:} La distancia geométrica entre números representables consecutivos no es constante. Se vuelve más densa cerca del cero y crece exponencialmente a medida que aumenta el exponente.
\end{itemize}

\subsection*{9. Redondeo y Truncado de Hardware}
Dado un valor exacto y un límite de mantisa de $m$ dígitos:
\begin{itemize}
    \item \textbf{Algoritmo de Truncado (Corte):} El hardware retiene los primeros $m$ dígitos de la mantisa y descarta sistemáticamente toda cifra posterior ($d_{m+1}$ en adelante), subestimando la magnitud final.
    \item \textbf{Algoritmo de Redondeo al más cercano:} El hardware evalúa $d_{m+1}$. Si $d_{m+1} < \beta/2$, trunca. Si $d_{m+1} \ge \beta/2$, incrementa una unidad a la última posición retenida $d_m$, propagando acarreo si es necesario.
\end{itemize}

\subsection*{10. Error Absoluto y Error Relativo}
Si $x$ es el valor real exacto y $fl(x)$ su representación de máquina:
\begin{itemize}
    \item \textbf{Error Absoluto ($E_{abs}$):} $E_{abs} = |x - fl(x)|$. Mide la magnitud de la distancia directa.
    \item \textbf{Error Relativo ($E_{rel}$):} $E_{rel} = \frac{|x - fl(x)|}{|x|}$ (para $x \neq 0$). Proporciona el desvío en proporción al valor original, haciéndolo invariante ante la escala numérica.
\end{itemize}

\subsection*{11. Épsilon de la Máquina ($\epsilon_{mach}$)}
Es el número positivo más pequeño tal que la computadora reconoce matemáticamente que $1 + \epsilon_{mach} > 1$. Representa la cota superior del error relativo del sistema.
\begin{itemize}
    \item Para sistemas que truncan: $\epsilon_{mach} = \beta^{1-m}$
    \item Para sistemas que redondean: $\epsilon_{mach} = \frac{1}{2} \beta^{1-m}$
\end{itemize}

\subsection*{12. Cancelación Catastrófica}
Pérdida severa de cifras significativas críticas que ocurre al restar dos aproximaciones de máquina casi idénticas. Esto elimina los dígitos principales exactos, desplazando los errores de redondeo o la "basura de memoria" hacia las posiciones más significativas.

\section*{Unidad 2 — Ecuaciones Diferenciales: Problemas de Valores Iniciales}

\subsection*{Introducción y Motivación Geométrica}
Un Problema de Valor Inicial (PVI) nos entrega una regla de crecimiento general dictada por una Ecuación Diferencial Ordinaria (EDO) de la forma $y' = f(t,y)$, junto con una única coordenada exacta de partida $y(t_0) = y_0$. 

El error numérico nace inevitablemente de la \textbf{discretización}. Mientras que la solución analítica $y(t)$ es una curva continua que ajusta su pendiente en cada instante infinitesimal, la computadora avanza mediante saltos discretos de tamaño $h$. Al "congelar" la pendiente durante todo el ancho del paso $h$, el algoritmo ignora la micro-curvatura real de la función, generando una divergencia geométrica en cada salto.

\subsection*{Anatomía de los Errores Numéricos}

En la resolución de EDOs, es vital no confundir el error que genera el algoritmo en un instante aislado con el error que se acumula al final de la simulación.

\section*{Problema de Valor Inicial y Existencia}

El problema de valor inicial para una ecuación diferencial de orden 1 se define mediante la ecuación $x'(t) = f(t, x(t))$ sujeta a una condición inicial $x(t_0) = a$[cite: 1].

\begin{teorema}[Teorema 8.2 - Existencia y Unicidad]
Si $f(t, x)$ es continua y Lipschitz en $x$, es decir existe $L > 0$ tal que $|f(t, x) - f(t, y)| \le L|x - y|$, para $t \in I$, con $I \subset \mathbb{R}$ intervalo y para todo $x \in \mathbb{R}$[cite: 1]. Entonces para cualquier valor $a \in \mathbb{R}$ existe una única función derivable $x(t)$ definida en $I$ que verifica $x'(t) = f(t, x(t))$ para todo $t \in I$ y $x(t_0) = a$[cite: 1].
\end{teorema}

\section*{Métodos de Un Paso}

Los métodos de un paso buscan generar $n$ valores $x_1, \dots, x_n$ que aproximen a $x(t_1), \dots, x(t_n)$ utilizando únicamente la información del paso inmediato anterior. 
Su forma matemática general es:
\begin{equation}
    x_{i+1} = x_i + h\Phi(t_i, x_i, h)
\end{equation}
donde $\Phi$ se conoce como la \textbf{función de incremento}.

\subsection*{13. Método de Euler}
Aproxima la derivada de la función $x$ mediante cocientes incrementales[cite: 1]. Dada una aproximación $x_i$ en $t_i$, se define la siguiente aproximación como:
\begin{equation}
    x_{i+1} = x_i + h f(t_i, x_i)
\end{equation}
Este método corresponde a la estructura general con $\Phi(t, x, h) = f(t, x)$[cite: 1].

\subsection*{14. Métodos de Taylor}
Generalizan el método de Euler utilizando desarrollos de Taylor de orden superior[cite: 1]. Si definimos el operador $T_k(t, x, h) = \sum_{j=1}^k \frac{h^{j-1}}{j!} D^j f(t, x)$, la fórmula de recurrencia de orden $k$ resulta en:
\begin{equation}
    x_{i+1} = x_i + h T_k(t_i, x_i, h)
\end{equation}
Su principal desventaja es que requieren encontrar y evaluar múltiples derivadas sucesivas totales de $f(t, x)$

\subsection*{15. Métodos de Runge-Kutta}
Tienen como objetivo obtener aproximaciones del mismo orden que los métodos de Taylor pero sin la necesidad computacional de calcular derivadas de $f$
\begin{itemize}
    \item \textbf{Euler Modificado (Orden 2):} $x_{i+1} = x_i + h f\left(t_i + \frac{h}{2}, x_i + \frac{h}{2}f(t_i, x_i)\right)$
    \item \textbf{Método de Heun (Orden 2):} $x_{i+1} = x_i + \frac{h}{2} [f(t_i, x_i) + f(t_{i+1}, x_i + h f(t_i, x_i))]$
\end{itemize}

\section*{Análisis del Error y Convergencia}

\subsection*{16. Error absoluto local y error de truncamiento local}

El Error Local Absoluto representa el desvío geométrico físico generado en un único salto aislado de tamaño $h$, asumiendo que el algoritmo arrancó ese instante desde una coordenada matemáticamente perfecta en la curva real. Esta distancia exacta conforma el término exacto que el método decide descartar y nace directamente de truncar la Serie de Taylor. Para visualizar su origen, se expande la solución analítica verdadera, la cual avanza sumando infinitos términos:
$$y(t_{n+1}) = y(t_n) + h y'(t_n) + \frac{h^2}{2!} y''(t_n) + \frac{h^3}{3!} y'''(t_n) + \dots$$Si aplicamos un algoritmo de primer orden como el Método de Euler Explícito, la computadora aproxima este avance trazando una línea recta, reteniendo únicamente la posición y la primera derivada:
$$y_{n+1} = y_n + h f(t_n, y_n)$$Al restar esta aproximación lineal de la curva analítica exacta, todo lo que el algoritmo tiró a la basura conforma el error del salto. Como el primer término descartado (el de grado 2) es el más grande, este domina el comportamiento del desvío:
$$\textbf{Error Absoluto Local} =\frac{h^2}{2} y''(\xi) = \mathcal{O}(h^2)$$

Llamaremos $\tau_i$ al \textbf{error de truncamiento local} que se comete en el paso $i$-ésimo. Este error está dado por la expresión:
\begin{equation}
    x(t_i + h) = x(t_i) + h\Phi(t_i, x(t_i), h) + h\tau_i
\end{equation}
Es decir, $\tau_i = \frac{x(t_i + h) - x(t_i)}{h} - \Phi(t_i, x(t_i), h)$

El error de truncamiento local, a nivel gráfico, es el error geométrico que se genera en un \textbf{único paso} iterativo $h$, asumiendo que el algoritmo arrancó desde una coordenada espacial perfecta. Representa la basura analítica que el método decide descartar.

\begin{proposicion}
El error de truncamiento local del método de Taylor de orden $k$ está dado por $\tau = \frac{h^k}{(k+1)!} x^{(k+1)}(\xi)$, para algún $\xi \in (t, t + h)$[cite: 1]. En consecuencia, este método es efectivamente de orden $k$.
\end{proposicion}

\begin{proposicion}
El error de truncamiento local de un método de Runge-Kutta de orden $k$ verifica $\tau = \mathcal{O}(h^k)$.
\end{proposicion}

En gran parte de la literatura, el error de truncado local se define como un \textit{residual por unidad de paso} (denotado como $\tau$). Esto se obtiene dividiendo el error absoluto geométrico por el tamaño del paso $h$.


$$\textbf{Error de Truncamiento Local} = \tau = \frac{\mathcal{O}(h^2)}{h} = \mathcal{O}(h) $$

\subsection*{17. Orden de un Método ($k$)}
Decir que un método numérico es de ``Orden $k$'' establece una garantía estricta sobre su precisión final a largo plazo. 
\begin{itemize}
    \item \textbf{Definición:} Diremos que el método es de orden $k$ si el error de truncamiento local satisface $\tau = \mathcal{O}(h^k)$
    \item \textbf{Implicancia:} Si reducimos el tamaño de paso $h$ a la mitad, el error global máximo garantizado se encogerá por un factor estricto de $2^k$.
    \item \textbf{Relación Local-Global:} Para lograr un Error Global de $\mathcal{O}(h^k)$, el error absoluto que el algoritmo descarta en cada salto local individual debe ser de un grado matemáticamente superior: $\mathcal{O}(h^{k+1})$.
\end{itemize}

\subsection*{18. Error global}
\begin{teorema}
Consideremos el método de un paso $x_{i+1} = x_i + h\Phi(t_i, x_i, h)$, con $\Phi$ una función Lipschitz en la segunda variable (es decir, existe $K$ tal que $|\Phi(t, x, h) - \Phi(t, y, h)| \le K|x - y|$). Entonces, se tiene:
\begin{equation}
    |x(T) - x_n| \le \frac{\tau_{max}}{K} \left(e^{K(T - t_0)} - 1\right)
\end{equation}
donde $\tau_{max} = \max_{1 \le i \le n} \{|\tau_i|\}$ con $\tau_i$ el error de truncamiento local del método en el paso $i$.
\end{teorema}

\begin{proposicion}[Proposición 8.11 - Condición de Convergencia]
Para un método de un paso con función de incremento $\Phi$ continua en sus tres variables, se tiene que $\lim_{h \to 0} \tau_{max} = 0$ si y sólo si $\Phi(t, x, 0) = f(t, x)$.
\end{proposicion}

\subsection*{Una aplicación del Teorema de error global (Acotación en Euler)}
Este teorema unifica el error local del algoritmo y el amplificador $L$ de la ecuación diferencial para construir un techo matemático inquebrantable. Para el Método de Euler en un intervalo $[t_0, t_f]$, la cota del error global garantiza que:
\begin{equation}
    |y(t_n) - y_n| \le \frac{M}{2L} \left( e^{L(t_n - t_0)} - 1 \right) h
\end{equation}
\textbf{Anatomía de la Cota:}
\begin{itemize}
    \item \textbf{$M = \max|y''(t)|$:} Representa el error máximo originado localmente por usar rectas en lugar de curvas.
    \item \textbf{$h$:} Evidencia que el Método de Euler es estrictamente de orden 1 ($\mathcal{O}(h)$).
    \item \textbf{$e^{L(\dots)}$:} El factor de amplificación. Dicta qué tan rápido explota la propagación del error a lo largo del tiempo de simulación debido a la sensibilidad $L$ de la ecuación original.
\end{itemize}


\subsection*{19. Sistemas de EDOs y formulación Matricial y Vectorial}
Un Problema de Valor Inicial para un sistema acoplado de $m$ ecuaciones diferenciales ordinarias de primer orden se formula formalmente mediante el siguiente sistema simultáneo:

\begin{equation}
\begin{cases}
    y_1'(t) = f_1(t, y_1(t), y_2(t), \dots, y_m(t)), & y_1(t_0) = y_{1,0} \\
    y_2'(t) = f_2(t, y_1(t), y_2(t), \dots, y_m(t)), & y_2(t_0) = y_{2,0} \\
    \quad \vdots & \quad \vdots \\
    y_m'(t) = f_m(t, y_1(t), y_2(t), \dots, y_m(t)), & y_m(t_0) = y_{m,0}
\end{cases}
\end{equation}


Este sistema de $m$ Ecuaciones Diferenciales Ordinarias (EDOs) de primer orden acopladas se puede representar analíticamente mediante notación vectorial como:
\begin{equation}
    \mathbf{Y}'(t) = \mathbf{F}(t, \textbf{Y}(t))
\end{equation}
Donde el vector de estado $\mathbf{Y}(t) = [y_1(t), y_2(t), \dots, y_m(t)]^T$ y la función vectorial $\mathbf{F}: \mathbb{R} \times \mathbb{R}^m \to \mathbb{R}^m$. Para un Problema de Valor Inicial (PVI), la información de partida se concentra en el vector constante $\mathbf{Y}(t_0) = \mathbf{Y}_0$.

Los algoritmos de un paso (como Euler o Runge-Kutta) mantienen su arquitectura algebraica, pero sus operaciones mutan hacia el álgebra lineal, evaluando cada componente del vector de forma simultánea:
\begin{equation}
    \mathbf{Y}_{n+1} = \mathbf{Y}_n + h \mathbf{\Phi}(t_n, \mathbf{Y}_n, h)
\end{equation}

\subsection*{20. La Matriz Jacobiana y el Límite de Sensibilidad}
En sistemas multidimensionales, la tasa de variación que define la constante amplificadora de errores muta de una derivada parcial escalar hacia la \textbf{Matriz Jacobiana} $J \in \mathbb{R}^{m \times m}$:
\begin{equation}
    J(t, \mathbf{Y}) = \frac{\partial \mathbf{F}}{\partial \mathbf{Y}} = 
    \begin{bmatrix}
        \frac{\partial f_1}{\partial y_1} & \cdots & \frac{\partial f_1}{\partial y_m} \\
        \vdots & \ddots & \vdots \\
        \frac{\partial f_m}{\partial y_1} & \cdots & \frac{\partial f_m}{\partial y_m}
    \end{bmatrix}
\end{equation}

\subsection*{21. Autovalores y Estabilidad Asintótica}
El comportamiento de las trayectorias alrededor del equilibrio está gobernado estrictamente por el signo de la parte real de los autovalores evaluados en la matriz J(t, \mathbf{Y}):
\begin{itemize}
    \item \textbf{Estabilidad Asintótica (Atractor):} Ocurre si \textbf{absolutamente todos} los autovalores poseen una parte real estrictamente negativa ($\text{Re}(\lambda_i) < 0$). Físicamente, el término exponencial de la solución analítica ($e^{\lambda t}$) exhibe un decaimiento incondicional. Cualquier perturbación matemática o desvío numérico se extingue asintóticamente hacia $0$ a medida que $t \to \infty$.
    \item \textbf{Inestabilidad (Repulsor):} Si existe al menos un autovalor con parte real positiva ($\text{Re}(\lambda_i) > 0$), las trayectorias vecinas divergirán exponencialmente alejándose del equilibrio.
    \item \textbf{Estabilidad Marginal:} La presencia de autovalores puramente nulos ($\lambda_i = 0$) indica que el sistema no diverge ni converge de manera estricta en esa dirección particular, revelando la existencia de un continuo de equilibrio.
\end{itemize}

Que el modelo físico sea analíticamente estable (decaimiento asegurado) no garantiza que el método numérico logre simularlo sin explotar.
\begin{itemize}
    \item \textbf{El Límite del Paso ($h$):} Para que un método numérico explícito (como Euler) logre imitar este decaimiento estable, el factor de amplificación discreto dictado por cada autovalor debe caer estrictamente dentro de su Región de Estabilidad Absoluta (por ejemplo, cumpliendo $|1 + h\lambda_i| < 1$).
    \item \textbf{Diagnóstico de Rigidez:} Si el sistema es asintóticamente estable pero la matriz Jacobiana revela autovalores con partes reales negativas de \textbf{gran magnitud} (por ejemplo, $\lambda = -5000$), el sistema se clasifica oficialmente como \textit{rígido} (stiff).
    \item \textbf{Transición Obligatoria a Métodos Implícitos:} La rigidez extrema fuerza a los esquemas explícitos a utilizar pasos temporales $h$ microscópicos para evitar colapsar numéricamente. Esta restricción asfixiante exige invariablemente abandonar la iteración directa y aplicar \textbf{esquemas implícitos} (A-estables), los cuales acoplan la incógnita futura y permiten avanzar con tamaños de paso masivos absorbiendo la rigidez sin perder estabilidad.
\end{itemize}

\subsection*{22. La constante Lipschitz en sistemas de ED0s}
La Constante de Lipschitz escalar $L$ se redefine utilizando una norma matricial subordinada evaluada sobre el dominio del problema: $L = \sup \|J(t, \mathbf{Y})\|$. Los autovalores del Jacobiano ($\lambda_i$) dictan la estabilidad del sistema; si poseen partes reales negativas de gran magnitud, el sistema se clasifica como \textit{rígido} (stiff), exigiendo invariablemente el uso de esquemas implícitos.

\subsection*{23. Reducción de Ecuaciones de Orden $n$ a Sistemas Vectoriales}

Cuando un modelo físico es gobernado por una única Ecuación Diferencial Ordinaria de orden $n$, de la forma:
\begin{equation}
    y^{(n)}(t) = f\left(t, y(t), y'(t), y''(t), \dots, y^{(n-1)}(t)\right)
\end{equation}
con $n$ condiciones auxiliares concentradas simultáneamente en el instante inicial $t_0$:
\begin{equation}
    y(t_0) = \alpha_1, \quad y'(t_0) = \alpha_2, \quad \dots, \quad y^{(n-1)}(t_0) = \alpha_n
\end{equation}

Es matemáticamente imposible aplicar algoritmos como Euler o Runge-Kutta de forma directa, ya que estos exigen un sistema de primer orden. Para resolverlo, el problema escalar se "vectoriza" introduciendo $n$ variables de estado auxiliares que encapsulan las derivadas sucesivas:

\begin{align*}
    u_1(t) &= y(t) \\
    u_2(t) &= y'(t) \\
    u_3(t) &= y''(t) \\
    &\vdots \\
    u_n(t) &= y^{(n-1)}(t)
\end{align*}

Al derivar cada una de estas variables respecto al tiempo, se encadenan formando un sistema acoplado de $n$ ecuaciones de primer orden:
\begin{equation}
\begin{cases}
    u_1'(t) = u_2(t) \\
    u_2'(t) = u_3(t) \\
    \quad \vdots \\
    u_{n-1}'(t) = u_n(t) \\
    u_n'(t) = f(t, u_1, u_2, \dots, u_n) \quad \leftarrow \text{(La EDO original)}
\end{cases}
\end{equation}

\textbf{Formulación Matricial Definitiva:}
Este acoplamiento colapsa en un vector de estado $\mathbf{U}(t)$ y un campo direccional $\mathbf{F}$, recuperando la arquitectura estándar $\mathbf{U}' = \mathbf{F}(t, \mathbf{U})$:
\begin{equation}
\mathbf{U}'(t) = 
\begin{pmatrix}
    u_1' \\ u_2' \\ \vdots \\ u_{n-1}' \\ u_n'
\end{pmatrix}
=
\begin{pmatrix}
    u_2 \\ u_3 \\ \vdots \\ u_n \\ f(t, u_1, u_2, \dots, u_n)
\end{pmatrix},
\quad
\mathbf{U}(t_0) = 
\begin{pmatrix}
    \alpha_1 \\ \alpha_2 \\ \vdots \\ \alpha_{n-1} \\ \alpha_n
\end{pmatrix}
\end{equation}
Al estar toda la información $\alpha_i$ definida en $t_0$, el problema mantiene el estatus de Problema de Valor Inicial y puede proyectarse hacia el futuro mediante álgebra lineal iterativa.

\vspace{0.5cm}
\subsection*{24. Reducción de PVI Escalares a Sistemas Vectoriales}

Para aplicar algoritmos de un paso (como Euler o Taylor) a Ecuaciones Diferenciales Ordinarias (EDO) de orden superior, es requisito matemático reducir el problema escalar a un sistema acoplado de primer orden mediante la introducción de variables de estado auxiliares.

\subsubsection*{Planteo del Problema Físico}
Sea un PVI de orden 2, gobernado por la EDO $y'' = f(t, y, y')$ y con toda su información auxiliar concentrada en el instante de partida $t_0$:
\begin{equation}
\begin{cases}
    y''(t) = -ky(t) - cy'(t) \quad \text{(Ej: Oscilador amortiguado)} \\
    y(t_0) = \alpha_1 \\
    y'(t_0) = \alpha_2
\end{cases}
\end{equation}

\subsubsection*{Vectorización del Estado (El Sistema Equivalente)}
Definimos variables de estado para la posición ($u_1 = y$) y la velocidad ($u_2 = y'$). Derivándolas respecto al tiempo, estructuramos un sistema acoplado que se empaqueta matricialmente:
\begin{equation}
\mathbf{U}'(t) = 
\begin{pmatrix}
    u_1' \\ u_2'
\end{pmatrix}
=
\begin{pmatrix}
    u_2 \\ -k u_1 - c u_2
\end{pmatrix}
= \mathbf{F}(t, \mathbf{U}), 
\quad \mathbf{U}(t_0) = \begin{pmatrix} \alpha_1 \\ \alpha_2 \end{pmatrix}
\end{equation}

\subsubsection*{Resolución Algorítmica con Taylor de Orden 2}
El Método de Taylor de Orden 2, aplicado sobre la arquitectura matricial, avanza la simulación evaluando simultáneamente todas las componentes vectoriales y reteniéndolas hasta la segunda derivada temporal:
\begin{equation}
    \mathbf{U}_{n+1} = \mathbf{U}_n + h \mathbf{U}_n' + \frac{h^2}{2} \mathbf{U}_n''
\end{equation}

Para calcular la matriz $\mathbf{U}_n''$, derivamos analíticamente el campo vectorial $\mathbf{F}$ respecto al tiempo aplicando regla de la cadena:
\begin{equation}
\mathbf{U}'' = \frac{d}{dt} \mathbf{F}(t, \mathbf{U}) = 
\begin{pmatrix}
    u_2' \\ -k u_1' - c u_2'
\end{pmatrix}
=
\begin{pmatrix}
    u_2' \\ -k u_2 - c u_2'
\end{pmatrix}
\end{equation}

Sustituyendo $\mathbf{U}'$ y $\mathbf{U}''$ evaluadas en el instante iterativo $n$, el avance discreto se opera evaluando:
\begin{equation}
\begin{pmatrix}
    u_{1, n+1} \\ u_{2, n+1}
\end{pmatrix}
=
\begin{pmatrix}
    u_{1, n} \\ u_{2, n}
\end{pmatrix}
+ h 
\begin{pmatrix}
    u_{2, n} \\ -k u_{1, n} - c u_{2, n}
\end{pmatrix}
+ \frac{h^2}{2}
\begin{pmatrix}
    -k u_{1, n} - c u_{2, n} \\ -k u_{2, n} - c(-k u_{1, n} - c u_{2, n})
\end{pmatrix}
\end{equation}

\subsubsection*{Error Local de Truncamiento Vectorial}
Al cortar la Serie de Taylor analítica en el término de grado 2, el algoritmo descarta los términos superiores, asumiendo un error local que hereda la estructura del vector columna[cite: 8]. Para este sistema acoplado, el \textbf{Error Local Absoluto} de la iteración está acotado por la norma máxima del primer término descartado (el resto de Lagrange de grado 3)[cite: 8, 10]:
\begin{equation}
    \|\mathbf{E}_{local}\| \le \frac{h^3}{3!} \max \|\mathbf{U}'''(t)\| = \mathcal{O}(h^3)
\end{equation}

\section*{Unidad 2 — Ecuaciones Diferenciales: Problemas de Valores de Contorno (PVC)}

A diferencia de los Problemas de Valor Inicial (PVI), donde toda la información se concentra en un único instante de tiempo permitiendo una marcha directa, los Problemas de Valores de Contorno (PVC) imponen condiciones en los extremos de un dominio físico cerrado $[a, b]$. Esta fractura de la información anula los algoritmos de marcha iterativa simples y exige resolver todo el dominio espaciado de manera simultánea.

\subsection*{25. Ecuaciones Ordinarias de Segundo Orden y Existencia}

La forma estándar de un PVC lineal unidimensional de segundo orden es:
\begin{equation}
    y''(x) = p(x)y'(x) + q(x)y(x) + r(x), \quad \forall x \in [a, b]
\end{equation}
Sujeto a las condiciones de frontera de Dirichlet: $y(a) = \alpha$ y $y(b) = \beta$.

Para que este problema tenga sentido físico y computacional, debe garantizarse que la solución no solo exista, sino que sea única (evitando que la curva oscile caóticamente entre los dos extremos).

\begin{teorema}[Existencia y Unicidad de PVC]
Supongamos que $p(x)$, $q(x)$ y $r(x)$ son funciones continuas en el intervalo $[a, b]$. Si además se cumple la restricción geométrica estricta de que:
\begin{equation}
    q(x) > 0 \quad \text{para todo } x \in [a, b]
\end{equation}
entonces el problema de valores de contorno tiene una \textbf{única solución} $y(x)$.
\end{teorema}
\textit{Nota analítica:} La condición $q(x) > 0$ le otorga ``rigidez'' a la curva matemática, prohibiéndole cruzar el eje o desarrollar múltiples ondulaciones que satisfagan los mismos bordes.


\section*{Discretización y el Método de Diferencias Finitas}

\subsection*{27. Discretización del dominio}

\begin{definition}
Sea $[a,b]$ un intervalo cerrado. Una \textbf{discretización} del dominio
consiste en reemplazar el dominio continuo por un conjunto finito de puntos,
denominados \textbf{nodos},
\[
\mathcal{M}_h=\{x_0,x_1,\ldots,x_N\},
\]
que aproximan el dominio original.
\end{definition}

En una malla uniforme, los nodos se encuentran igualmente espaciados. El
\textbf{paso de malla} $h$ se define mediante
\[
\boxed{
h=\frac{b-a}{N}
}
\tag{24}
\]
y los nodos quedan dados por
\[
\boxed{
x_i=a+ih,
\qquad i=0,1,\ldots,N.
}
\]

En particular,
\[
x_0=a,
\qquad
x_N=b.
\]

La función continua $y(x)$ queda representada sobre la malla mediante los
valores discretos
\[
\boxed{
y_i\approx y(x_i).
}
\]

\subsection*{28. Diferencias finitas}

\begin{definition}
Una \textbf{diferencia finita} es una aproximación algebraica de una
derivada, construida a partir de los valores de una función en nodos
discretos de una malla.

Estas aproximaciones se obtienen a partir del desarrollo de Taylor de la
función alrededor de los nodos de la malla.
\end{definition}

En una malla uniforme, los desarrollos de Taylor alrededor del nodo $x_i$
son
\[
y(x_i+h)
=
y(x_i)
+h\,y'(x_i)
+\frac{h^2}{2!}y''(x_i)
+\frac{h^3}{3!}y^{(3)}(x_i)
+\mathcal{O}(h^4),
\]
y
\[
y(x_i-h)
=
y(x_i)
-h\,y'(x_i)
+\frac{h^2}{2!}y''(x_i)
-\frac{h^3}{3!}y^{(3)}(x_i)
+\mathcal{O}(h^4).
\]

A partir de estas expresiones se obtienen los principales esquemas de
diferencias finitas.

\subsubsection*{Diferencia hacia adelante (Forward)}

La aproximación de la primera derivada mediante diferencia hacia adelante es
\[
\boxed{
y'(x_i)\approx
\frac{y_{i+1}-y_i}{h}.
}
\]

En particular,
\[
\frac{y_{i+1}-y_i}{h}
=
y'(x_i)+\mathcal{O}(h),
\]
por lo que el error de truncamiento asintótico es
\[
\boxed{
\mathcal{O}(h).
}
\]

\subsubsection*{Diferencia hacia atrás (Backward)}

La aproximación de la primera derivada mediante diferencia hacia atrás es
\[
\boxed{
y'(x_i)\approx
\frac{y_i-y_{i-1}}{h}.
}
\]

Análogamente,
\[
\frac{y_i-y_{i-1}}{h}
=
y'(x_i)+\mathcal{O}(h),
\]
por lo que el error de truncamiento asintótico es
\[
\boxed{
\mathcal{O}(h).
}
\]

\subsubsection*{PVC: diferencia centrada}

Para resolver numéricamente el PVC, se discretiza el dominio $[a,b]$ en
$N$ subintervalos de igual longitud mediante el paso de malla
\[
h=\frac{b-a}{N}.
\]

Se obtiene así el conjunto de nodos
\[
x_i=a+ih,
\qquad i=0,1,\ldots,N.
\]

Para aproximar la primera derivada, restamos los desarrollos de Taylor
alrededor de $x_i$, obteniendo
\[
y(x_i+h)-y(x_i-h)
=
2h\,y'(x_i)+\mathcal{O}(h^3).
\]

Por tanto,
\[
\boxed{
y'(x_i)\approx
\frac{y_{i+1}-y_{i-1}}{2h},
}
\tag{25}
\]
con error de truncamiento
\[
\boxed{
\mathcal{O}(h^2).
}
\]

Para la segunda derivada, sumamos los desarrollos de Taylor y obtenemos
\[
y(x_i+h)+y(x_i-h)
=
2y(x_i)+h^2y''(x_i)+\mathcal{O}(h^4).
\]

De aquí se sigue que
\[
\boxed{
y''(x_i)\approx
\frac{y_{i+1}-2y_i+y_{i-1}}{h^2},
}
\tag{26}
\]
con error de truncamiento
\[
\boxed{
\mathcal{O}(h^2).
}
\]

En consecuencia, las derivadas analíticas continuas se reemplazan por
aproximaciones algebraicas definidas sobre la malla discreta.

\section*{Estructuración Analítica de Problemas de Valores de Contorno (PVC)}

\subsection*{28. Discretización Espacial y Agrupación de Coeficientes}
Partimos de la Ecuación Diferencial Ordinaria (EDO) genérica lineal de segundo orden:
\begin{equation}
    y''(t) = u(t) + v(t)y(t) + w(t)y'(t)
\end{equation}
Aplicando el esquema de \textbf{diferencias finitas centradas de segundo orden} sobre una malla discreta de paso $h$, aproximamos las derivadas en el nodo interior $i$:
\begin{equation}
    \frac{y_{i+1} - 2y_i + y_{i-1}}{h^2} = u_i + v_i y_i + w_i \frac{y_{i+1} - y_{i-1}}{2h}
\end{equation}

Para construir el sistema lineal, multiplicamos toda la expresión por $h^2$ para eliminar denominadores y agrupamos estrictamente las variables incógnitas ($y_{i-1}, y_i, y_{i+1}$) a la izquierda y las constantes a la derecha. Expandiendo y reordenando:
\begin{align*}
    y_{i+1} - 2y_i + y_{i-1} &= h^2 u_i + h^2 v_i y_i + \frac{h w_i}{2} y_{i+1} - \frac{h w_i}{2} y_{i-1} \\
    \underbrace{\left(1 + \frac{h w_i}{2}\right)}_{a_{i-1}} y_{i-1} + \underbrace{\left(-2 - h^2 v_i\right)}_{d_i} y_i + \underbrace{\left(1 - \frac{h w_i}{2}\right)}_{c_i} y_{i+1} &= \underbrace{h^2 u_i}_{b_i}
\end{align*}
Esta agrupación define algebraicamente los coeficientes exactos de la matriz tridiagonal para el sistema iterativo $AY = B$.

\subsection*{29. Absorción de Fronteras (Condiciones de Dirichlet) y sistema matricial}
Para un dominio discretizado en $m$ nodos interiores, las fronteras físicas de los extremos imponen valores conocidos $y_0 = \alpha$ y $y_{m+1} = \beta$. Dado que no son incógnitas, se despejan hacia la derecha, siendo absorbidas por el vector constante de términos independientes $B$.

\begin{itemize}
    \item \textbf{Primer nodo interior ($i=1$):} La condición de frontera izquierda $\alpha$ impacta contra la subdiagonal del sistema.
    \begin{equation}
        d_1 y_1 + c_1 y_2 = b_1 - a_0 \alpha
    \end{equation}
    \item \textbf{Último nodo interior ($i=m$):} La condición de frontera derecha $\beta$ impacta contra la superdiagonal.
    \begin{equation}
        a_{m-1} y_{m-1} + d_m y_m = b_m - c_m \beta
    \end{equation}
\end{itemize}

El ensamblaje de todas las ecuaciones nodales conforma el sistema matricial tridiagonal completo AY=b. Aquí se hace explícito el mapeo exacto de los apuntes manuscritos a la estructura de banda:
\begin{equation}
\begin{pmatrix}
    d_1 & c_1 & 0 & \cdots & 0 \\
    a_1 & d_2 & c_2 & \cdots & 0 \\
    0 & a_2 & d_3 & \cdots & 0 \\
    \vdots & \ddots & \ddots & \ddots & \vdots \\
    0 & \cdots & 0 & a_{m-1} & d_m
\end{pmatrix}
\begin{pmatrix}
    y_1 \\
    y_2 \\
    y_3 \\
    \vdots \\
    y_m
\end{pmatrix}
=
\begin{pmatrix}
    b_1 - a_0 \alpha \\
    b_2 \\
    b_3 \\
    \vdots \\
    b_m - c_m \beta
\end{pmatrix}
\end{equation}
En particular, a la ahora de resolver este sistema, analíticamente, resulta complicado obtener el vector Y de manera explícita invirtiendo a la matriz

\subsection*{30. La Matriz diagonal Dominante}
Una matriz se considera estrictamente diagonal dominante si, en cada fila, el valor absoluto del elemento en la diagonal principal es estrictamente mayor que la suma de los valores absolutos del resto de los elementos de esa fila:
$|A_{ii}| > \sum_{j \neq i} |A_{ij}|$.

\begin{lema} [Lema de Hadamard]
Sea A una matriz diagonal dominante entonces A es no singular (inversible)
\end{lema} 

\begin{lema}[Herencia de la Dominancia Diagonal Estricta]
Sea $\mathbf{A} \in \mathbb{R}^{n \times n}$ una matriz estrictamente diagonal dominante (EDD). Entonces, toda \textbf{submatriz principal} $\mathbf{M}$ de tamaño $m \times m$ (con $1 \le m \le n$) hereda la propiedad de ser estrictamente diagonal dominante.
\end{lema}

\begin{proposicion} [Estabilidad Numérica del PVC]
Sea un Problema de Valores de Contorno (PVC) que satisface las hipótesis del Teorema de Existencia y Unicidad (con $q(x) > 0$). Si el paso de discretización $h$ se elige tal que:
\[ h < \frac{2}{\max |p(x)|} \]
entonces la matriz de coeficientes tridiagonal $\mathbf{A} \in \mathbb{R}^{n \times n}$ asociada al sistema discreto es \textbf{estrictamente diagonal dominante} (EDD).
\begin{itemize}
    \item {Por el Lema de Hadamard, la matriz es no singular (inversible). Esto asegura de forma absoluta que el sistema lineal discreto admite una solución única}
    \item {El sistema puede resolverse mediante Factorización LU optimizada o métodos iterativos (como Jacobi o Gauss-Seidel) garantizando la convergencia absoluta sin necesidad de pivoteo.}
\end{itemize}
\end{proposicion} 

Para garantizar que la matriz $A$ sea numéricamente estable y permita la convergencia sin necesidad de pivoteo al invertirla, demostramos formalmente que el valor absoluto de la diagonal principal supera a la suma de los valores absolutos del resto de la fila: $|d_i| > |a_{i-1}| + |c_i|$.

 Exigimos estructuralmente que el tamaño de paso $h$ sea lo suficientemente pequeño para garantizar la inecuación estricta $-1 \le \frac{w_i h}{2} \le 1$. Esto asegura matemáticamente que los términos $\left(1 - \frac{w_i h}{2}\right)$ y $\left(1 + \frac{w_i h}{2}\right)$ sean siempre positivos o nulos, permitiéndonos retirar las barras de valor absoluto sin alterar los signos.

Sumando las diagonales exteriores bajo esta hipótesis:
\begin{equation}
    |a_{i-1}| + |c_i| = \left|1 + \frac{w_i h}{2}\right| + \left|1 - \frac{w_i h}{2}\right| = \left(1 + \frac{w_i h}{2}\right) + \left(1 - \frac{w_i h}{2}\right) = 2
\end{equation}
Evaluando la diagonal principal (asumiendo por el teorema de unicidad que $v_i > 0$):
\begin{equation}
    |d_i| = |-2 - h^2 v_i| = 2 + h^2 v_i
\end{equation}
Dado que $v_i > 0$, la constante sumada garantiza una inecuación inquebrantable, concluyendo la demostración:
\begin{equation}
    2 + h^2 v_i > 2 \implies |d_i| > |a_{i-1}| + |c_i|
\end{equation}

\subsection*{31. Error Local de Truncamiento ($T_j$) en PVC}

Al discretizar una ecuación diferencial ordinaria de segundo orden lineal (un PVC) de la forma genérica $x''(t) = u(t) + v(t)x(t) + w(t)x'(t)$ mediante Diferencias Finitas Centradas, sustituimos los operadores analíticos por aproximaciones algebraicas.

Si evaluamos el esquema estrictamente sobre la \textbf{solución analítica exacta} $x(t_j)$ en un nodo discreto $j$, la igualdad deja de ser perfecta debido a los términos que el algoritmo desecha de la Serie de Taylor. Este residual numérico se define formalmente como el \textbf{Error Local de Truncamiento}, denotado por $T_j$.

Basado en la estructura algorítmica, el residual $T_j$ se define como:
\begin{equation}
    T_j = \left( \frac{x(t_{j+1}) - 2x(t_j) + x(t_{j-1})}{h^2} \right) - \left( u_j + v_j x(t_j) + w_j \frac{x(t_{j+1}) - x(t_{j-1})}{2h} \right)
\end{equation}

\begin{proposicion} [Condición de suavidad: $x \in C^4{([a,b])}$]               
Para acotar este residual $T_j$, necesitamos analizar los términos descartados por el método en cuestión que utilicemos. Para Taylor de órden 2, la aproximación centrada de la segunda derivada $x''(t_j)$ omite términos proporcionales a la cuarta derivada de la solución, $x^{(4)}(\xi)$. 
Por lo tanto, se exige como hipótesis estricta que la solución exacta pertenezca al espacio $C^4([a,b])$ (es decir, que admita hasta cuatro derivadas continuas). Bajo esta condición geométrica de suavidad, la cuarta derivada está acotada por un máximo en el intervalo, lo que garantiza inquebrantablemente que el error inyectado en cada nodo aislado escala de forma cuadrática respecto al paso $h$:
\begin{equation}
    \max_{1 \le j \le n-1} |T_j| \le \mathcal{O}(h^2)
\end{equation}
\end{proposicion}

\subsection*{32. Error Global y propagación}
Definimos el Error Global $e_i = x(t_i) - y_i$, midiendo la diferencia entre la solución analítica exacta evaluada en el nodo y la solución discreta calculada. Al insertar la solución analítica en nuestro esquema de diferencias finitas, el truncamiento de la Serie de Taylor inyecta un error residual local $E_i$.

Al restar el sistema numérico del sistema analítico, observamos que el vector de errores hereda exactamente la misma arquitectura matricial tridiagonal:
\begin{equation}
    a_{i-1} e_{i-1} + d_i e_i + c_i e_{i+1} = h^2 E_i
\end{equation}

De forma análoga al sistema original, el sistema de propagación de errores se estructura explícitamente en notación matricial:
\begin{equation}
\begin{pmatrix}
    d_1 & c_1 & 0 & \cdots & 0 \\
    a_1 & d_2 & c_2 & \cdots & 0 \\
    0 & a_2 & d_3 & \cdots & 0 \\
    \vdots & \ddots & \ddots & \ddots & \vdots \\
    0 & \cdots & 0 & a_{m-1} & d_m
\end{pmatrix}
\begin{pmatrix}
    e_1 \\
    e_2 \\
    e_3 \\
    \vdots \\
    e_m
\end{pmatrix}
= h^2
\begin{pmatrix}
    E_1 \\
    E_2 \\
    E_3 \\
    \vdots \\
    E_m
\end{pmatrix}
\end{equation}

La estabilidad de este sistema de errores está garantizada por la holgura de la dominancia diagonal que demostramos anteriormente:
\begin{equation}
    |d_i| - (|a_{i-1}| + |c_i|) = (2 + h^2 v_i) - 2 = h^2 v_i
\end{equation}
Esta "holgura" estrictamente positiva de valor $h^2 v_i$ actúa como la barrera estructural matemática que absorbe las perturbaciones, impidiendo que los errores de truncamiento locales $E_i$ se amplifiquen explosivamente al resolver el sistema matricial, garantizando la convergencia del esquema.

\vspace{0.5cm}

\subsection*{33. Convergencia y Error Global ($e_j$)}

El algoritmo real no opera sobre la solución exacta $x(t_j)$, sino que calcula vectores aproximados $y_j$ resolviendo un sistema de ecuaciones lineales de la forma $\mathbf{A}Y = \mathbf{b}$. La diferencia absoluta entre la realidad analítica y la simulación discreta en cada nodo es el \textbf{Error Global}, definido como $e_j = x(t_j) - y_j$.

\begin{teorema}[Convergencia Asintótica del PVC]
Supongamos que $x(t)$ es la solución analítica exacta del problema continuo de valores de contorno continuo, y satisface la hipótesis de suavidad $x \in C^4([a,b])$. Sea $y \in \mathbb{R}^n$ el vector numérico obtenido al resolver el sistema algebraico discreto por diferencias finitas centradas.

Si el sistema continuo cumple el teorema de unicidad y la matriz resultante $\mathbf{A}$ es estrictamente diagonal dominante, entonces existe una constante estructural $C > 0$ (independiente del paso de discretización $h$) tal que la norma del máximo del Error Global satisface:
\begin{equation}
    \max_{1 \le j \le n} |e_j| \le C \cdot h^2
\end{equation}
\end{teorema}

\subsection*{Explicación Operativa del Teorema}
Este teorema demuestra que el esquema numérico matricial es \textbf{convergente}. 

A diferencia de los Problemas de Valor Inicial (PVI), donde el error local crece de forma exponencial debido al efecto amplificador del avance iterativo en el tiempo, en los Problemas de Valores de Contorno el sistema se resuelve en su totalidad de manera simultánea. 
Gracias a las propiedades estructurales de la matriz tridiagonal dominante $\mathbf{A}$ (cuya inversa está estrictamente acotada, $\|\mathbf{A}^{-1}\| \le M$), los errores de truncamiento locales $T_j$ inyectados en la ecuación no se amplifican descontroladamente.

Como resultado, la cota del Error Global ($e_j$) "hereda" directamente el mismo orden de convergencia \textit{orden} del Error Local de Truncamiento ($T_j$). Es decir, los dos son de orden $\mathcal{O}(h^2)$, certificando que si decides refinar tu malla espacial a la mitad, la precisión total de tu simulación del sistema físico se cuadruplicará de forma garantizada.

\section*{Introducción a las Ecuaciones en Derivadas Parciales (EDPs)}


\vspace{0.3cm}

Al escalar el modelado a múltiples dimensiones (espacio y tiempo simultáneamente), abandonamos las EDOs y entramos en el dominio de las Ecuaciones en Derivadas Parciales (EDP). 

\subsection*{34. Forma general de un sistema de EDP}
Para que una EDP evolutiva posea una solución matemática única y determinista, la ecuación gobernante debe estar encapsulada por condiciones auxiliares que restrinjan todo su dominio. Tomando como ejemplo la \textbf{Ecuación del Calor} unidimensional, el sistema completo se estructura bajo la siguiente llave fundamental:

\begin{equation}
\begin{cases}
    u_t(x,t) = \alpha \cdot u_{xx}(x,t) & \text{en } 0 < x < L, \quad t > 0 \quad \text{\textbf{(Ecuación Gobernante)}} \\
    u(x,0) = f(x) & \text{en } 0 \le x \le L \quad \text{\textbf{(Condición Inicial)}} \\
    u(0,t) = g_1(t) & \text{para } t > 0 \quad \text{\textbf{(Condición de Borde - Izquierda)}} \\
    u(L,t) = g_2(t) & \text{para } t > 0 \quad \text{\textbf{(Condición de Borde - Derecha)}}
\end{cases}
\end{equation}

\textbf{Desglose de la Información:}

\begin{itemize}
    \item \textbf{Condición Inicial:} Define el estado del sistema en el instante de arranque temporal $t = 0$. Ejemplo: $u(x, 0) = f(x)$. Responde a la pregunta: \textit{¿Cómo estaba la distribución antes de que encendiera el cronómetro?}
    \item \textbf{Condición de Borde (o Frontera):} Restringe el comportamiento de la función en los límites espaciales geométricos durante todo el tiempo de simulación. 
    \begin{itemize}
        \item \textit{Condición de Dirichlet:} Fija un valor exacto (Ej. bordes a $0^\circ$C).
        \item \textit{Condición de Neumann:} Fija la derivada o flujo (Ej. bordes perfectamente aislados, flujo de calor cero).
    \end{itemize}
\end{itemize}

\subsection*{35. Grillas Computacionales y Notación de Índices}
Para atacar numéricamente una EDP bidimensional (espacio $x$ y tiempo $t$), se construye una malla o \textbf{grilla espacio-temporal}:
\begin{itemize}
    \item Discretizamos el espacio con pasos $\Delta x$: \quad $x_j = j \Delta x$
    \item Discretizamos el tiempo con pasos $\Delta t$: \quad $t_n = n \Delta t$
\end{itemize}
\textbf{La Notación Estándar:} La aproximación numérica de la función continua $u(x_j, t_n)$ se escribe utilizando un subíndice para el nodo espacial y un superíndice para el paso iterativo temporal:
\begin{equation}
    U_j^n \approx u(x_j, t_n)
\end{equation}

Para representar que estamos parados en el nodo espacial $j$ durante la iteración temporal $n$, la notación matricial estándar ubica al espacio en el subíndice y al tiempo en el superíndice: $U_j^n \approx u(x_j, t_n)$.

\vspace{0.5cm}
\begin{center}
\begin{tikzpicture}[scale=1.5]
    % Ejes principales
    \draw[->, thick] (0,0) -- (5,0) node[right] {$x$ (Espacio)};
    \draw[->, thick] (0,0) -- (0,3.5) node[above] {$t$ (Tiempo)};
    
    % Grilla interna (lineas punteadas)
    \draw[step=1cm, gray, very thin, dashed] (0,0) grid (4.2, 3.2);
    
    % Dibujo de los nodos
    \foreach \x in {0,1,2,3,4} {
        \foreach \y in {0,1,2,3} {
            \fill (\x,\y) circle (2pt);
        }
    }
    
    % Etiquetas del eje X (Espacio)
    \node[below] at (0,0) {$x_0=0$};
    \node[below] at (1,0) {$x_1$};
    \node[below] at (2,0) {$x_2$};
    \node[below] at (3,0) {$x_3$};
    \node[below] at (4,0) {$x_N=L$};
    
    % Etiquetas del eje Y (Tiempo)
    \node[left] at (0,1) {$t_1$};
    \node[left] at (0,2) {$t_2$};
    \node[left] at (0,3) {$t_3$};
    
    % Resaltado de un nodo genérico
    \fill[red] (2,2) circle (3pt) node[above right, text=red, font=\bfseries] {$U_j^n \approx u(x_j, t_n)$};
    
    % Resaltado de las Condiciones Auxiliares
    \draw[ultra thick, blue] (0,0) -- (4,0) node[midway, below=0.6cm, align=center] {\textbf{Condición Inicial} \\ $u(x,0) = f(x)$};
    \draw[ultra thick, green!60!black] (0,0) -- (0,3.2) node[midway, left=0.8cm, align=center] {\textbf{Borde Izquierdo} \\ $u(0,t)$};
    \draw[ultra thick, green!60!black] (4,0) -- (4,3.2) node[midway, right=0.3cm, align=center] {\textbf{Borde Derecho} \\ $u(L,t)$};
\end{tikzpicture}
\end{center}
\vspace{0.5cm}

La grilla revela el inmenso costo algebraico de las EDPs: para avanzar desde la capa temporal $t_1$ hacia la capa temporal futura $t_2$, el algoritmo debe resolver simultáneamente todos los nodos de la fila horizontal evaluando sus acoplamientos espaciales. Esto se ejecuta resolviendo iterativamente sistemas de la forma $\mathbf{A} \vec{U}^{n+1} = \vec{U}^n + \vec{b}$.

\subsection*{Esquemas explícitos e implícitos}

En la resolución numérica de ecuaciones en derivadas parciales dependientes del tiempo, una de las decisiones fundamentales consiste en determinar cómo se discretiza la derivada temporal. Esta elección conduce, de manera natural, a la distinción entre \emph{esquemas explícitos} e \emph{implícitos}.

Considérese una ecuación evolutiva escrita de forma abstracta como
\begin{equation}
    \frac{\partial u}{\partial t}
    =
    \mathcal{L}(u),
\end{equation}
donde $\mathcal{L}$ representa un operador espacial, que puede contener derivadas espaciales de distintos órdenes y, eventualmente, términos no lineales.

Se introduce una discretización temporal uniforme
\begin{equation}
    t^n = n\Delta t,
    \qquad
    n=0,1,2,\ldots,
\end{equation}
donde $\Delta t>0$ es el paso temporal. Se denota por
\begin{equation}
    u^n \approx u(\cdot,t^n)
\end{equation}
la aproximación numérica de la solución en el instante $t^n$.

La diferencia esencial entre los métodos explícitos e implícitos reside en el instante temporal en el que se evalúa el operador espacial $\mathcal{L}$.

\subsection*{Esquemas explícitos}

En un esquema explícito, la solución en el instante futuro $t^{n+1}$ se calcula exclusivamente a partir de información correspondiente al instante actual $t^n$ y a instantes anteriores.

De forma general, un esquema explícito de un paso puede escribirse como
\begin{equation}
    u^{n+1}
    =
    \Phi(u^n),
\end{equation}
donde el operador $\Phi$ depende únicamente de cantidades ya conocidas.

La principal ventaja de este tipo de métodos es que el avance temporal se realiza directamente, sin necesidad de resolver sistemas algebraicos en cada paso. Sin embargo, esta simplicidad computacional suele venir acompañada de restricciones de estabilidad sobre el tamaño del paso temporal $\Delta t$.

\subsubsection*{Método de Euler explícito}

El método de Euler explícito se obtiene aproximando la derivada temporal mediante una diferencia progresiva:
\begin{equation}
    \frac{\partial u}{\partial t}(t^n)
    \approx
    \frac{u^{n+1}-u^n}{\Delta t}.
\end{equation}

Evaluando el operador espacial en el instante conocido $t^n$, se obtiene
\begin{equation}
    \frac{u^{n+1}-u^n}{\Delta t}
    =
    \mathcal{L}(u^n).
\end{equation}

Por tanto,
\begin{equation}
    \boxed{
    u^{n+1}
    =
    u^n
    +
    \Delta t\,\mathcal{L}(u^n)
    }
\end{equation}

Esta expresión muestra claramente el carácter explícito del método: para calcular $u^{n+1}$ solamente se necesitan valores ya conocidos en el instante $t^n$.

El método de Euler explícito es de primer orden en el tiempo. En efecto, mediante el desarrollo de Taylor
\begin{equation}
    u(t^{n+1})
    =
    u(t^n)
    +
    \Delta t\,u_t(t^n)
    +
    \frac{\Delta t^2}{2}u_{tt}(\xi),
\end{equation}
para algún $\xi\in(t^n,t^{n+1})$, se observa que el error local de truncamiento es de orden
\begin{equation}
    \mathcal{O}(\Delta t^2),
\end{equation}
mientras que el error global acumulado es
\begin{equation}
    \mathcal{O}(\Delta t).
\end{equation}

\subsubsection*{Ejemplo: ecuación del calor}

Considérese la ecuación unidimensional del calor
\begin{equation}
    \frac{\partial u}{\partial t}
    =
    \alpha
    \frac{\partial^2u}{\partial x^2},
    \qquad
    \alpha>0.
\end{equation}

Sea una malla espacial uniforme
\begin{equation}
    x_i=i\Delta x,
\end{equation}
y denótese por
\begin{equation}
    u_i^n
    \approx
    u(x_i,t^n).
\end{equation}

La derivada segunda espacial se aproxima mediante diferencias centradas:
\begin{equation}
    \frac{\partial^2u}{\partial x^2}(x_i,t^n)
    \approx
    \frac{
        u_{i+1}^n
        -
        2u_i^n
        +
        u_{i-1}^n
    }{\Delta x^2}.
\end{equation}

Aplicando Euler explícito,
\begin{equation}
    \frac{u_i^{n+1}-u_i^n}{\Delta t}
    =
    \alpha
    \frac{
        u_{i+1}^n
        -
        2u_i^n
        +
        u_{i-1}^n
    }{\Delta x^2}.
\end{equation}

Definiendo
\begin{equation}
    r
    =
    \frac{\alpha\Delta t}{\Delta x^2},
\end{equation}
se obtiene
\begin{equation}
    \boxed{
    u_i^{n+1}
    =
    r u_{i-1}^n
    +
    (1-2r)u_i^n
    +
    r u_{i+1}^n
    }
\end{equation}

La expresión anterior permite calcular directamente cada valor $u_i^{n+1}$ a partir de los valores de la capa temporal anterior.

Para este problema, el método es estable únicamente bajo una restricción del tipo
\begin{equation}
    r
    =
    \frac{\alpha\Delta t}{\Delta x^2}
    \leq
    \frac{1}{2}.
\end{equation}

Equivalentemente,
\begin{equation}
    \boxed{
    \Delta t
    \leq
    \frac{\Delta x^2}{2\alpha}
    }
\end{equation}

Esta condición muestra una característica importante de los métodos explícitos: al refinar la malla espacial, es necesario reducir también el paso temporal para conservar la estabilidad.

En particular, si $\Delta x$ se divide por dos, entonces $\Delta t$ debe reducirse aproximadamente por un factor cuatro.

\subsection*{Esquemas implícitos}

En un esquema implícito, la ecuación que determina la solución en el instante $t^{n+1}$ contiene también valores desconocidos correspondientes a dicho instante.

De forma abstracta, un esquema implícito puede escribirse como
\begin{equation}
    F(u^{n+1},u^n)=0.
\end{equation}

Por tanto, el cálculo de $u^{n+1}$ no se obtiene mediante una fórmula directa, sino que requiere resolver una ecuación algebraica o, después de discretizar el espacio, un sistema de ecuaciones algebraicas.

La principal desventaja de los métodos implícitos es, por tanto, su mayor coste computacional por paso temporal. A cambio, suelen presentar propiedades de estabilidad considerablemente mejores que los métodos explícitos.

\subsubsection*{Método de Euler implícito}

El método de Euler implícito aproxima la derivada temporal mediante
\begin{equation}
    \frac{\partial u}{\partial t}(t^{n+1})
    \approx
    \frac{u^{n+1}-u^n}{\Delta t},
\end{equation}
pero evalúa el operador espacial en el nuevo instante temporal:
\begin{equation}
    \frac{u^{n+1}-u^n}{\Delta t}
    =
    \mathcal{L}(u^{n+1}).
\end{equation}

Por tanto,
\begin{equation}
    \boxed{
    u^{n+1}
    =
    u^n
    +
    \Delta t\,
    \mathcal{L}(u^{n+1})
    }
\end{equation}

A diferencia del método explícito, $u^{n+1}$ aparece en ambos lados de la ecuación.

El método de Euler implícito es también de primer orden en el tiempo:
\begin{equation}
    e_{\mathrm{global}}
    =
    \mathcal{O}(\Delta t).
\end{equation}

Sin embargo, sus propiedades de estabilidad son generalmente superiores a las de Euler explícito.

\subsubsection*{Ejemplo: ecuación del calor}

Aplicando Euler implícito a
\begin{equation}
    u_t=\alpha u_{xx},
\end{equation}
se obtiene
\begin{equation}
    \frac{u_i^{n+1}-u_i^n}{\Delta t}
    =
    \alpha
    \frac{
        u_{i+1}^{n+1}
        -
        2u_i^{n+1}
        +
        u_{i-1}^{n+1}
    }{\Delta x^2}.
\end{equation}

Introduciendo nuevamente
\begin{equation}
    r
    =
    \frac{\alpha\Delta t}{\Delta x^2},
\end{equation}
resulta
\begin{equation}
    u_i^{n+1}
    -
    u_i^n
    =
    r
    \left(
        u_{i+1}^{n+1}
        -
        2u_i^{n+1}
        +
        u_{i-1}^{n+1}
    \right).
\end{equation}

Reordenando los términos,
\begin{equation}
    \boxed{
    -r u_{i-1}^{n+1}
    +
    (1+2r)u_i^{n+1}
    -
    r u_{i+1}^{n+1}
    =
    u_i^n
    }
\end{equation}

El valor $u_i^{n+1}$ depende, por tanto, de los valores igualmente desconocidos
\begin{equation}
    u_{i-1}^{n+1}
    \qquad\text{y}\qquad
    u_{i+1}^{n+1}.
\end{equation}

Consecuentemente, todos los valores de la nueva capa temporal deben determinarse simultáneamente.

Si se agrupan las incógnitas interiores en el vector
\begin{equation}
    U^{n+1}
    =
    \begin{pmatrix}
        u_1^{n+1}\\
        u_2^{n+1}\\
        \vdots\\
        u_{N-1}^{n+1}
    \end{pmatrix},
\end{equation}
el esquema puede escribirse matricialmente como
\begin{equation}
    A U^{n+1}
    =
    b^n,
\end{equation}
donde, para condiciones de contorno adecuadas, la matriz $A$ presenta la estructura tridiagonal
\begin{equation}
    A
    =
    \begin{pmatrix}
        1+2r & -r     &        &        & 0\\
        -r   & 1+2r   & -r     &        & \\
             & \ddots & \ddots & \ddots & \\
             &        & -r     & 1+2r   & -r\\
        0    &        &        & -r     & 1+2r
    \end{pmatrix}.
\end{equation}

Así, en cada paso temporal debe resolverse
\begin{equation}
    \boxed{
    A U^{n+1}
    =
    b^n
    }
\end{equation}

Es importante observar que, aunque formalmente podría escribirse
\begin{equation}
    U^{n+1}
    =
    A^{-1}b^n,
\end{equation}
en la práctica numérica no se calcula explícitamente la matriz inversa. En su lugar, se resuelve directamente el sistema lineal mediante métodos apropiados, por ejemplo eliminación gaussiana, factorización LU o, en el caso tridiagonal, el algoritmo de Thomas.

Para la ecuación del calor, Euler implícito es incondicionalmente estable desde el punto de vista de la estabilidad lineal. En particular, no requiere una condición del tipo
\begin{equation}
    \Delta t
    \leq
    C\Delta x^2
\end{equation}
para evitar inestabilidades numéricas.

Esto no significa, sin embargo, que $\Delta t$ pueda elegirse arbitrariamente grande. La estabilidad y la precisión son conceptos distintos: un paso temporal excesivamente grande puede producir una solución estable pero poco precisa.

\subsection*{Comparación entre Euler explícito y Euler implícito}

Las diferencias esenciales entre ambos métodos pueden resumirse de la siguiente manera:
\begin{center}
\begin{tabular}{c|c|c}
    & Euler explícito & Euler implícito\\
    \hline
    Operador espacial
    &
    $\mathcal{L}(u^n)$
    &
    $\mathcal{L}(u^{n+1})$
    \\[0.2cm]
    Cálculo de $u^{n+1}$
    &
    Directo
    &
    Mediante sistema algebraico
    \\[0.2cm]
    Orden temporal
    &
    $\mathcal{O}(\Delta t)$
    &
    $\mathcal{O}(\Delta t)$
    \\[0.2cm]
    Coste por paso
    &
    Bajo
    &
    Mayor
    \\[0.2cm]
    Estabilidad
    &
    Condicional en numerosos problemas
    &
    Generalmente superior
\end{tabular}
\end{center}

La diferencia conceptual puede expresarse de forma sencilla:

\begin{quote}
En un método explícito, el estado futuro se construye utilizando únicamente información conocida del presente. En un método implícito, el propio estado futuro forma parte de la ecuación que debe satisfacer.
\end{quote}

\subsection*{Método de Crank--Nicolson}

El método de Crank--Nicolson constituye una alternativa intermedia entre los métodos de Euler explícito e implícito. Su idea fundamental consiste en aproximar el operador espacial mediante el promedio de sus valores en los instantes $t^n$ y $t^{n+1}$.

Partiendo de
\begin{equation}
    u_t
    =
    \mathcal{L}(u),
\end{equation}
se propone
\begin{equation}
    \frac{u^{n+1}-u^n}{\Delta t}
    =
    \frac{1}{2}
    \left[
        \mathcal{L}(u^n)
        +
        \mathcal{L}(u^{n+1})
    \right].
\end{equation}

Por tanto,
\begin{equation}
    \boxed{
    u^{n+1}
    =
    u^n
    +
    \frac{\Delta t}{2}
    \left[
        \mathcal{L}(u^n)
        +
        \mathcal{L}(u^{n+1})
    \right]
    }
\end{equation}

Aunque utiliza información de los dos niveles temporales, Crank--Nicolson es un método implícito, puesto que $u^{n+1}$ aparece dentro del operador espacial.

Una de sus principales ventajas respecto de los métodos de Euler es que posee orden dos en el tiempo:
\begin{equation}
    e_{\mathrm{global}}
    =
    \mathcal{O}(\Delta t^2).
\end{equation}

Por tanto, para pasos temporales suficientemente pequeños, Crank--Nicolson suele proporcionar una aproximación temporal significativamente más precisa que Euler explícito o Euler implícito.

\subsubsection*{Crank--Nicolson para la ecuación del calor}

Para
\begin{equation}
    u_t
    =
    \alpha u_{xx},
\end{equation}
Crank--Nicolson conduce a
\begin{equation}
    \frac{u_i^{n+1}-u_i^n}{\Delta t}
    =
    \frac{\alpha}{2}
    \left[
        \frac{
            u_{i+1}^{n}
            -
            2u_i^{n}
            +
            u_{i-1}^{n}
        }{\Delta x^2}
        +
        \frac{
            u_{i+1}^{n+1}
            -
            2u_i^{n+1}
            +
            u_{i-1}^{n+1}
        }{\Delta x^2}
    \right].
\end{equation}

Definiendo
\begin{equation}
    r
    =
    \frac{\alpha\Delta t}{\Delta x^2},
\end{equation}
se obtiene
\begin{equation}
    u_i^{n+1}-u_i^n
    =
    \frac{r}{2}
    \left[
        u_{i+1}^n
        -
        2u_i^n
        +
        u_{i-1}^n
        +
        u_{i+1}^{n+1}
        -
        2u_i^{n+1}
        +
        u_{i-1}^{n+1}
    \right].
\end{equation}

Agrupando en el lado izquierdo las incógnitas del instante $n+1$,
\begin{equation}
    -\frac{r}{2}u_{i-1}^{n+1}
    +
    (1+r)u_i^{n+1}
    -
    \frac{r}{2}u_{i+1}^{n+1}
    =
    \frac{r}{2}u_{i-1}^{n}
    +
    (1-r)u_i^{n}
    +
    \frac{r}{2}u_{i+1}^{n}.
\end{equation}

Así,
\begin{equation}
    \boxed{
    -\frac{r}{2}u_{i-1}^{n+1}
    +
    (1+r)u_i^{n+1}
    -
    \frac{r}{2}u_{i+1}^{n+1}
    =
    \frac{r}{2}u_{i-1}^{n}
    +
    (1-r)u_i^{n}
    +
    \frac{r}{2}u_{i+1}^{n}
    }
\end{equation}

En forma matricial, el esquema puede expresarse como
\begin{equation}
    A U^{n+1}
    =
    B U^n
    +
    g^n,
\end{equation}
donde $g^n$ contiene, en su caso, las contribuciones de las condiciones de contorno.

Las matrices correspondientes al problema interior tienen la forma
\begin{equation}
    A
    =
    \begin{pmatrix}
        1+r & -r/2 &       &       & 0\\
        -r/2 & 1+r & -r/2 &       & \\
             & \ddots & \ddots & \ddots & \\
             &       & -r/2 & 1+r & -r/2\\
        0    &       &       & -r/2 & 1+r
    \end{pmatrix},
\end{equation}
y
\begin{equation}
    B
    =
    \begin{pmatrix}
        1-r & r/2 &       &       & 0\\
        r/2 & 1-r & r/2   &       & \\
            & \ddots & \ddots & \ddots & \\
            &       & r/2 & 1-r & r/2\\
        0   &       &       & r/2 & 1-r
    \end{pmatrix}.
\end{equation}

Por tanto, en cada instante temporal se resuelve
\begin{equation}
    \boxed{
    A U^{n+1}
    =
    B U^n
    +
    g^n
    }
\end{equation}

Al igual que Euler implícito, Crank--Nicolson presenta una gran estabilidad para problemas parabólicos lineales. Para la ecuación del calor es incondicionalmente estable en el sentido de von Neumann.

Sin embargo, Crank--Nicolson no posee el mismo comportamiento disipativo que Euler implícito. Para pasos temporales muy grandes pueden aparecer oscilaciones numéricas asociadas a modos de alta frecuencia, aun cuando el esquema permanezca estable. Por este motivo, la estabilidad incondicional no debe interpretarse como una garantía de buena precisión para cualquier valor de $\Delta t$.

\subsection*{Interpretación mediante el método $\theta$}

Los tres esquemas anteriores pueden reunirse dentro de una única familia conocida como \emph{método $\theta$}:
\begin{equation}
    \frac{u^{n+1}-u^n}{\Delta t}
    =
    (1-\theta)\mathcal{L}(u^n)
    +
    \theta\mathcal{L}(u^{n+1}),
\end{equation}
donde
\begin{equation}
    0\leq\theta\leq1.
\end{equation}

Los casos principales son:
\begin{equation}
    \boxed{
    \begin{aligned}
        \theta=0
        &\quad\Longrightarrow\quad
        \text{Euler explícito},
        \\
        \theta=1
        &\quad\Longrightarrow\quad
        \text{Euler implícito},
        \\
        \theta=\frac12
        &\quad\Longrightarrow\quad
        \text{Crank--Nicolson}.
    \end{aligned}
    }
\end{equation}

Esta formulación permite interpretar los distintos métodos según el peso asignado a la información de cada nivel temporal.

Cuando $\theta=0$, toda la información espacial procede del instante conocido $t^n$, por lo que el esquema es completamente explícito.

Cuando $\theta=1$, toda la información espacial se evalúa en $t^{n+1}$, de modo que el esquema es completamente implícito.

Finalmente, cuando
\begin{equation}
    \theta=\frac12,
\end{equation}
se utiliza un promedio entre ambos niveles temporales, dando lugar al método de Crank--Nicolson.

\subsection*{Esquemas lineales y no lineales}

Conviene distinguir entre el carácter implícito del método y la naturaleza del sistema algebraico que debe resolverse.

Si el operador $\mathcal{L}$ es lineal, la discretización implícita conduce habitualmente a un sistema lineal de la forma
\begin{equation}
    A U^{n+1}
    =
    b^n.
\end{equation}

En cambio, si la EDP es no lineal, puede aparecer un sistema no lineal
\begin{equation}
    F(U^{n+1})=0.
\end{equation}

En tal caso, cada paso temporal puede requerir el uso de un procedimiento iterativo, como el método de Newton.

Por ejemplo, para una ecuación esquemática
\begin{equation}
    u_t
    =
    \mathcal{N}(u),
\end{equation}
con $\mathcal{N}$ no lineal, Euler implícito proporciona
\begin{equation}
    u^{n+1}
    -
    \Delta t\,\mathcal{N}(u^{n+1})
    -
    u^n
    =
    0.
\end{equation}

La incógnita $u^{n+1}$ debe obtenerse entonces resolviendo una ecuación no lineal.

Por tanto, el término \emph{implícito} no significa simplemente ``resolver una matriz'', sino, más generalmente, que la solución desconocida del nuevo instante aparece dentro de la ecuación que la define.

\subsection*{Estabilidad frente a precisión}

Una distinción fundamental en la elección de un esquema temporal es la diferencia entre \emph{estabilidad} y \emph{precisión}.

La estabilidad estudia si los errores numéricos se amplifican o permanecen controlados durante la evolución temporal. La precisión, en cambio, mide qué tan próxima se encuentra la solución numérica de la solución exacta.

Un método implícito puede ser estable para valores grandes de $\Delta t$, pero ello no implica necesariamente que la solución obtenida sea precisa.

Por ejemplo, Euler implícito aplicado a la ecuación del calor permite utilizar valores de $\Delta t$ mucho mayores que Euler explícito desde el punto de vista de la estabilidad. Sin embargo, dado que su error temporal global es
\begin{equation}
    \mathcal{O}(\Delta t),
\end{equation}
un valor excesivamente grande de $\Delta t$ puede producir una aproximación muy pobre.

Análogamente, Crank--Nicolson presenta un error
\begin{equation}
    \mathcal{O}(\Delta t^2),
\end{equation}
pero un paso temporal excesivamente grande puede seguir siendo inadecuado desde el punto de vista de la resolución temporal de la solución.

Por tanto,
\begin{equation}
    \boxed{
    \text{estabilidad}
    \;\not\Rightarrow\;
    \text{precisión}
    }
\end{equation}

\subsection*{Las ventajas y desventajas de los métodos explícitos e implícitos}

La elección entre un método explícito e implícito representa un compromiso entre coste computacional, estabilidad y precisión.

El método de Euler explícito presenta la forma
\begin{equation}
    \boxed{
    u^{n+1}
    =
    u^n
    +
    \Delta t\,
    \mathcal{L}(u^n)
    }
\end{equation}
y permite avanzar directamente en el tiempo, aunque suele estar sometido a restricciones de estabilidad.

El método de Euler implícito viene dado por
\begin{equation}
    \boxed{
    u^{n+1}
    =
    u^n
    +
    \Delta t\,
    \mathcal{L}(u^{n+1})
    }
\end{equation}
y requiere resolver un sistema algebraico en cada paso temporal, pero presenta mejores propiedades de estabilidad.

Finalmente, el método de Crank--Nicolson se escribe como
\begin{equation}
    \boxed{
    \frac{u^{n+1}-u^n}{\Delta t}
    =
    \frac{1}{2}
    \left[
        \mathcal{L}(u^n)
        +
        \mathcal{L}(u^{n+1})
    \right]
    }
\end{equation}
y combina información de los instantes $t^n$ y $t^{n+1}$. Es un método implícito de segundo orden temporal y constituye una de las discretizaciones temporales más utilizadas para problemas parabólicos.

Desde un punto de vista conceptual, puede resumirse la diferencia fundamental de la siguiente manera:
\begin{equation}
    \boxed{
    \begin{array}{ccl}
        \text{Explícito}
        &:&
        \text{el futuro se calcula a partir del presente},
        \\[0.2cm]
        \text{Implícito}
        &:&
        \text{el futuro se determina resolviendo una ecuación que lo contiene}.
    \end{array}
    }
\end{equation}

\subsection*{Estabilidad y el Principio del Máximo Discreto}
En la simulación iterativa de EDPs sobre miles de nodos, el peligro crítico es la \textbf{inestabilidad numérica}: la amplificación caótica de los errores de redondeo de la máquina a medida que avanzamos en el eje del tiempo.

\begin{teorema}[Estabilidad en Norma del Máximo]
Sea $u_j^n$ la solución exacta que brindaría el esquema en diferencias finitias (si la computadora tuviera precisión infinita) y $v_j^n$ la solución numérica real calculada (que arrastra la "basura" de redondeo del hardware). El error absoluto aislado que existe en el nodo espacial $j$ durante la iteración temporal $n$ es:
\begin{equation}
    e_j^n = |u_j^n - v_j^n|
\end{equation}

Se define la \textbf{Norma del Máximo} del vector de errores de toda esa capa temporal como la magnitud de la mayor perturbación existente en la malla:
\begin{equation}
    E^n = \max_j |e_j^n| = \|\vec{e}^{\,n}\|_\infty
\end{equation}

El esquema de discretización se clasifica como \textbf{estrictamente estable} si cumple el Principio del Máximo Discreto. Matemáticamente, exige que exista una constante $C \le 1$ tal que la cota máxima del error no se amplifique jamás al avanzar un paso de tiempo:
\begin{equation}
    E^{n+1} \le C \cdot E^n \implies \|\vec{e}^{\,n+1}\|_\infty \le \|\vec{e}^{\,n}\|_\infty
\end{equation}
Esta inecuación inquebrantable garantiza que cualquier perturbación o desvío introducido en el algoritmo se disipará o se mantendrá acotado, evitando que la simulación explote iterativamente hacia el infinito.
\end{teorema}

Al escalar la complejidad a más de una dimensión (por ejemplo, evaluar la temperatura de una placa de metal en función del espacio $(x,y)$ y del tiempo $t$), abandonamos las EDOs y entramos en el dominio de las \textbf{EDPs}. 


\subsection*{36. Operación Matricial en las EDP}
Al sustituir las derivadas parciales por diferencias finitas (por ejemplo, con el esquema implícito Backward Euler para el tiempo y diferencias centradas para el espacio), el estado futuro de toda la malla espacial en el tiempo $n+1$ queda matemáticamente ligado. 
La marcha temporal abandona la idea de proyectar un solo punto (como en EDOs) y se convierte en una \textbf{sucesión de sistemas de ecuaciones lineales}. Cada iteración temporal exige resolver un sistema matricial de la forma:
\begin{equation}
    \mathbf{A} \vec{U}^{n+1} = \vec{U}^n + \vec{b}
\end{equation}
Donde $\vec{U}^n$ es todo el estado espacial actual y $\vec{U}^{n+1}$ es todo el estado espacial futuro.

\subsection*{37. Estabilidad y el Principio del Máximo Discreto}
En las EDPs, el error puede acumularse y explotar violentamente si la discretización no es estable. El análisis de estabilidad busca demostrar que las perturbaciones o errores no crecen descontroladamente al avanzar el tiempo $n$.

\begin{teorema}[Cota de Error de la Norma del Máximo]
Sea $u_j^n$ la solución exacta del esquema en diferencias discretas y $v_j^n$ la solución numérica calculada (que arrastra errores de redondeo de la máquina). Definimos el error en el nodo $j$ durante el instante $t_n$ como $e_j^n = |u_j^n - v_j^n|$. 
\end{teorema}

\section*{Unidad 4 — Norma y Condición de una Matriz}

\subsection*{38. Normas Matriciales}

En el estudio del álgebra lineal numérica, es fundamental cuantificar el ``tamaño'' de los vectores y de las transformaciones lineales (matrices). Para ello, extendemos el concepto de norma vectorial al espacio de las matrices. 

\begin{definicion}[Norma Matricial Inducida o Reducida]
Sea $A \in \mathbb{R}^{n \times m}$ una matriz y consideremos una norma vectorial $\|\cdot\| : \mathbb{R}^n \to \mathbb{R}_{\ge 0}$. Se define la \textit{norma matricial inducida} (o norma reducida) por la norma vectorial como el máximo factor de estiramiento que la matriz puede aplicarle a un vector:
\begin{equation}
    \|A\| = \max_{x \ne 0} \frac{\|Ax\|}{\|x\|}
\end{equation}
\end{definicion}

De esta definición se desprenden propiedades estructurales muy importantes. Por la propia construcción del máximo, se garantiza que para cualquier vector $x$ se cumple la cota $\|Ax\| \le \|A\| \|x\|$. Además, la norma inducida es submultiplicativa, lo que significa que dados los productos de matrices conformables, la norma del producto no excede el producto de las normas: $\|AB\| \le \|A\| \|B\|$.

Existen también normas matriciales que no provienen directamente de inducir una norma vectorial, como es el caso de la \textbf{Norma de Frobenius}, definida como la raíz cuadrada de la suma de todos los elementos de la matriz al cuadrado: $\|A\|_F = \sqrt{\sum_{i,j} |a_{ij}|^2}$.



\subsection*{39. Relación con las Normas Vectoriales Clásicas}
Recordemos las normas vectoriales más utilizadas:
\begin{itemize}
    \item Norma 1: $\|x\|_1 = \sum_{i=1}^n |x_i|$
    \item Norma 2 (Euclídea): $\|x\|_2 = \sqrt{\sum_{i=1}^n |x_i|^2}$
    \item Norma Infinito: $\|x\|_\infty = \max_{1 \le i \le n} |x_i|$
\end{itemize}

Cuando inducimos una norma matricial a partir de la norma infinito vectorial, obtenemos que $\|A\|_\infty$ equivale al máximo de la suma de los valores absolutos de las filas de la matriz:
\begin{equation}
    \|A\|_\infty = \max_{1 \le i \le n} \sum_{j=1}^m |a_{ij}|
\end{equation}

\begin{ejemplo}
Consideremos la matriz $A = \begin{pmatrix} 1 & -1 & 4 \\ 3 & 1 & 3 \end{pmatrix}$. Para calcular su norma infinito, evaluamos la suma de los módulos en cada fila:
\begin{itemize}
    \item Fila 1: $|1| + |-1| + |4| = 6$
    \item Fila 2: $|3| + |1| + |3| = 7$
\end{itemize}
Como el valor máximo ocurre en la segunda fila, concluimos que $\|A\|_\infty = 7$.
\end{ejemplo}

\paragraph{Normas matriciales inducidas clásicas.}

Para una matriz $A=(a_{ij})\in\R^{n\times n}$ se tienen las expresiones

\[
    \boxed{
    \|A\|_1
    =
    \max_{1\le j\le n}
    \sum_{i=1}^n |a_{ij}|
    }
\]

y

\[
    \boxed{
    \|A\|_\infty
    =
    \max_{1\le i\le n}
    \sum_{j=1}^n |a_{ij}|
    }.
\]

Es decir, la norma $1$ matricial corresponde a la máxima suma
absoluta por columnas, mientras que la norma infinito corresponde
a la máxima suma absoluta por filas.

Para la norma $2$ inducida se tiene

\[
    \boxed{
    \|A\|_2
    =
    \sqrt{\rho(A^TA)}
    =
    \sigma_{\max}(A),
    }
\]

donde $\sigma_{\max}(A)$ es el mayor valor singular de $A$.

En particular, si $A$ es simétrica,

\[
    \|A\|_2
    =
    \rho(A)
    =
    \max_i |\lambda_i(A)|.
\]

\begin{proposicion}[Comparación entre normas matriciales]
Sea $A\in\R^{n\times n}$. Entonces

\[
    \frac{1}{\sqrt n}\|A\|_\infty
    \leq
    \|A\|_2
    \leq
    \sqrt n\,\|A\|_\infty.
\]

Análogamente,

\[
    \frac{1}{\sqrt n}\|A\|_1
    \leq
    \|A\|_2
    \leq
    \sqrt n\,\|A\|_1.
\]
\end{proposicion}

Estas desigualdades permiten trasladar estimaciones obtenidas en una
norma a otra, siempre teniendo presente que las constantes dependen
de la dimensión $n$.

\subsection*{Equivalencia de Normas en Dimensión Finita}

\begin{teorema}[Equivalencia de Normas]
Sea $V$ un espacio vectorial de dimensión finita y sean 
$\|\cdot\|_a$ y $\|\cdot\|_b$ dos normas definidas sobre $V$.
Entonces existen constantes $c,C>0$, independientes de $x$, tales que
\[
c\|x\|_a
\leq
\|x\|_b
\leq
C\|x\|_a,
\qquad \forall x\in V.
\]
En consecuencia, todas las normas sobre un espacio vectorial de dimensión finita son equivalentes.
\end{teorema}

La equivalencia de normas no significa que ambas normas asignen el mismo valor a cada elemento, sino que sus valores pueden compararse mediante constantes multiplicativas independientes del elemento considerado.

Por ejemplo, en $\mathbb{R}^n$, las normas
\[
\|x\|_1,
\qquad
\|x\|_2,
\qquad
\|x\|_\infty
\]
no son, en general, iguales, pero sí son equivalentes.

En particular, se cumplen las desigualdades
\[
\|x\|_\infty
\leq
\|x\|_2
\leq
\|x\|_1,
\]
y además
\[
\|x\|_1
\leq
n\|x\|_\infty,
\qquad
\|x\|_2
\leq
\sqrt{n}\,\|x\|_\infty.
\]

\subsection*{40. Radio Espectral y Matrices Simétricas}

Un concepto íntimamente ligado al comportamiento asintótico y a las normas de las matrices es el radio espectral. Para ello, recordemos primero el Teorema Espectral.

\begin{teorema}[Teorema Espectral (Versión general)]
Sea $A \in \mathbb{C}^{n \times n}$ es una matriz hermitiana, entonces existe una base ortonormal de autovectores de A
\end{teorema}

Obs: Para el caso de $A \in \mathbb{R}^{n \times n}$ (acá hablamos de matrices simétricas), se cumple este teorema y además garantiza que los autovalores son reales          

\begin{definicion}[Radio Espectral]
El radio espectral de una matriz cuadrada $A$, denotado como $\rho(A)$, se define como el máximo de los módulos de sus autovalores:
\begin{equation}
    \rho(A) = \max \{|\lambda_i| : \lambda_i \text{ es autovalor de } A\}
\end{equation}
Es decir, si $\lambda_{max}$ es el autovalor de mayor magnitud, entonces $\rho(A) = |\lambda_{max}|$.
\end{definicion}

Para una matriz genérica, el radio espectral es siempre una cota inferior de cualquier norma inducida ($\rho(A) \le \|A\|$). Sin embargo, cuando la matriz posee una estructura geométrica particular, como la simetría, esta relación se convierte en una igualdad exacta para la norma 2.

\begin{teorema}[Norma 2 de Matrices Simétricas]
Si $A \in \mathbb{R}^{n \times n}$ es una matriz simétrica, entonces su norma 2 coincide exactamente con su radio espectral:
\begin{equation}
    \|A\|_2 = \rho(A)
\end{equation}
\end{teorema}

\begin{proof}
Por el Teorema Espectral, al ser $A$ simétrica, sabemos que posee una Base Ortonormal (BON) conformada íntegramente por sus autovectores $\{v_1, \dots, v_n\}$. Para estos autovectores se cumple que $Av_i = \lambda_i v_i$.

Dado cualquier vector $x \in \mathbb{R}^n$, podemos expresarlo como una combinación lineal de esta base: $x = \sum_{i=1}^n \alpha_i v_i$. 
Calculamos primero la norma al cuadrado del vector $x$:
\begin{equation}
    \|x\|_2^2 = \langle x, x \rangle = \left\langle \sum \alpha_i v_i, \sum \alpha_j v_j \right\rangle = \sum_{i=1}^n \alpha_i^2
\end{equation}
(los términos cruzados se anulan porque la base es ortonormal, es decir, $\langle v_i, v_j \rangle = \delta_{ij}$).

Ahora evaluamos la acción de la matriz sobre el vector: $Ax = A(\sum \alpha_i v_i) = \sum \alpha_i \lambda_i v_i$. 
Calculamos la norma al cuadrado de este nuevo vector:
\begin{equation}
    \|Ax\|_2^2 = \left\langle \sum \alpha_i \lambda_i v_i, \sum \alpha_j \lambda_j v_j \right\rangle = \sum_{i=1}^n \alpha_i^2 \lambda_i^2
\end{equation}
Sabiendo que para todo $i$, $\lambda_i^2 \le \lambda_{max}^2 = \rho(A)^2$, podemos acotar la suma:
\begin{equation}
    \|Ax\|_2^2 \le \lambda_{max}^2 \sum_{i=1}^n \alpha_i^2 = \rho(A)^2 \|x\|_2^2
\end{equation}
Tomando raíz cuadrada y dividiendo por $\|x\|_2$ (para $x \ne 0$), obtenemos que $\frac{\|Ax\|_2}{\|x\|_2} \le \rho(A)$, lo que prueba que $\|A\|_2 \le \rho(A)$.

Para probar que efectivamente se alcanza la igualdad, basta con evaluar el cociente en el autovector particular correspondiente al autovalor dominante. Sea $x = v_{max}$ tal que $A v_{max} = \lambda_{max} v_{max}$ (con $v_{max} \ne 0$). Entonces:
\begin{equation}
    \|A v_{max}\|_2 = \|\lambda_{max} v_{max}\|_2 = |\lambda_{max}| \|v_{max}\|_2 = \rho(A) \|v_{max}\|_2
\end{equation}
Al existir una dirección donde el estiramiento es exactamente $\rho(A)$, concluimos formalmente que $\|A\|_2 = \rho(A)$.
\end{proof}
\subsection*{41. Condicionamiento de una Matriz}

Cuando se resuelve un sistema lineal

\[
    Ax=b,
\]

es importante distinguir entre el cálculo de una solución y la
\emph{sensibilidad} del propio problema frente a perturbaciones en
los datos. Esta sensibilidad se cuantifica mediante el número de
condición.

\begin{definicion}[Número de Condición]
Sea $A\in\R^{n\times n}$ una matriz invertible y sea
$\|\cdot\|$ una norma matricial inducida. Se define

\[
    \boxed{
    \cond(A)
    =
    \|A\|\,\|A^{-1}\|
    }.
\]

Cuando se desea explicitar la norma empleada se escribe

\[
    \cond_p(A)
    =
    \|A\|_p\,\|A^{-1}\|_p.
\]

Si $A$ es singular, se adopta

\[
    \cond(A)=+\infty.
\]
\end{definicion}

El número de condición mide la sensibilidad relativa de la solución
de un sistema lineal frente a pequeñas perturbaciones en sus datos.

Una matriz con $\cond(A)$ cercano a $1$ se denomina
\emph{bien condicionada}; si $\cond(A)$ es grande, se dice que es
\emph{mal condicionada}.

\begin{proposicion}[Propiedades Básicas del Condicionamiento]
Sea $A$ invertible. Entonces:

\begin{enumerate}
    \item
    \[
        \boxed{\cond(A)\geq 1.}
    \]

    \item
    \[
        \boxed{\cond(A^{-1})=\cond(A).}
    \]

    \item Para todo $\alpha\neq0$,
    \[
        \boxed{\cond(\alpha A)=\cond(A).}
    \]

    \item Si $A$ y $B$ son invertibles,
    \[
        \boxed{
        \cond(AB)
        \leq
        \cond(A)\cond(B).
        }
    \]
\end{enumerate}
\end{proposicion}

\begin{proof}
Para la primera propiedad,

\[
    1
    =
    \|I\|
    =
    \|AA^{-1}\|
    \leq
    \|A\|\,\|A^{-1}\|
    =
    \cond(A).
\]

Además,

\[
    \cond(A^{-1})
    =
    \|A^{-1}\|\,\|A\|
    =
    \cond(A).
\]

Para $\alpha\neq0$,

\[
    \cond(\alpha A)
    =
    \|\alpha A\|
    \|(\alpha A)^{-1}\|
    =
    |\alpha|\|A\|
    \frac{1}{|\alpha|}\|A^{-1}\|
    =
    \cond(A).
\]

Finalmente,

\[
\begin{aligned}
    \cond(AB)
    &=
    \|AB\|\,\|(AB)^{-1}\| \\
    &=
    \|AB\|\,\|B^{-1}A^{-1}\| \\
    &\leq
    \|A\|\|B\|
    \|B^{-1}\|\|A^{-1}\| \\
    &=
    \cond(A)\cond(B).
\end{aligned}
\]
\end{proof}


\subsubsection*{Condicionamiento en norma 2}

Como

\[
    \|A\|_2=\sigma_{\max}(A),
\]

y, para $A$ invertible,

\[
    \|A^{-1}\|_2
    =
    \frac{1}{\sigma_{\min}(A)},
\]

se obtiene

\[
    \boxed{
    \cond_2(A)
    =
    \frac{\sigma_{\max}(A)}
         {\sigma_{\min}(A)}.
    }
\]

Esta expresión proporciona una interpretación geométrica del
condicionamiento: el número de condición compara el máximo y el mínimo
factor de estiramiento producido por la transformación lineal $A$.

Si $A$ es simétrica e invertible, sus valores singulares son los
módulos de sus autovalores, por lo que

\[
    \boxed{
    \cond_2(A)
    =
    \frac{\max_i|\lambda_i(A)|}
         {\min_i|\lambda_i(A)|}.
    }
\]

En particular, si $A$ es simétrica definida positiva,

\[
    \boxed{
    \cond_2(A)
    =
    \frac{\lambda_{\max}(A)}
         {\lambda_{\min}(A)}.
    }
\]

\begin{observacion}
Si $Q$ es ortogonal, entonces

\[
    Q^{-1}=Q^T,
\]

y todos sus valores singulares son iguales a $1$. Por tanto,

\[
    \boxed{\cond_2(Q)=1.}
\]

Esta propiedad será especialmente útil al estudiar la factorización
QR.
\end{observacion}


\subsubsection*{Comparación entre números de condición}

De las desigualdades entre normas matriciales se deduce que

\[
    \frac{1}{\sqrt n}\|A\|_\infty
    \leq
    \|A\|_2
    \leq
    \sqrt n\,\|A\|_\infty,
\]

y aplicando las mismas desigualdades a $A^{-1}$,

\[
    \boxed{
    \frac{1}{n}\cond_\infty(A)
    \leq
    \cond_2(A)
    \leq
    n\,\cond_\infty(A).
    }
\]

De forma análoga,

\[
    \boxed{
    \frac{1}{n}\cond_1(A)
    \leq
    \cond_2(A)
    \leq
    n\,\cond_1(A).
    }
\]

Estas cotas son especialmente útiles cuando resulta difícil calcular
directamente $\cond_2(A)$ pero se dispone de una expresión sencilla
para $\cond_1(A)$ o $\cond_\infty(A)$.


\subsubsection*{Equivalencia del condicionamiento en dimensión fija}

El espacio $\R^{n\times n}$ tiene dimensión finita $n^2$. Por lo tanto,
cualesquiera dos normas matriciales $\|\cdot\|_a$ y $\|\cdot\|_b$
son equivalentes: existen constantes $c,C>0$, independientes de $A$,
tales que

\[
    c\|A\|_a
    \leq
    \|A\|_b
    \leq
    C\|A\|_a.
\]

Aplicando estas desigualdades tanto a $A$ como a $A^{-1}$,

\[
    c^2\cond_a(A)
    \leq
    \cond_b(A)
    \leq
    C^2\cond_a(A).
\]

Por tanto, para una familia $A_\varepsilon$ de matrices de dimensión
fija,

\[
    \boxed{
    \cond_a(A_\varepsilon)\to\infty
    \iff
    \cond_b(A_\varepsilon)\to\infty.
    }
\]

\begin{observacion}
La hipótesis de dimensión fija es esencial. Si la dimensión $n$ varía,
las constantes de equivalencia pueden depender de $n$, por lo que no
se puede ignorar dicha dependencia en un argumento asintótico.
\end{observacion}


\subsection*{42. Teoremas de Perturbación}

El número de condición aparece de manera natural al estudiar cómo una
perturbación de los datos de un sistema lineal modifica su solución.


\subsubsection*{Perturbación del término independiente}

Consideremos

\[
    Ax=b
\]

y el sistema perturbado

\[
    A(x+\Delta x)=b+\Delta b.
\]

Restando ambas ecuaciones,

\[
    A\Delta x=\Delta b,
\]

y por tanto

\[
    \Delta x=A^{-1}\Delta b.
\]

\begin{teorema}[Perturbación del término independiente]
Sea $A$ invertible y $b\neq0$. Entonces

\[
    \boxed{
    \frac{1}{\cond(A)}
    \frac{\|\Delta b\|}{\|b\|}
    \leq
    \frac{\|\Delta x\|}{\|x\|}
    \leq
    \cond(A)
    \frac{\|\Delta b\|}{\|b\|}.
    }
\]
\end{teorema}

\begin{proof}
Como

\[
    \Delta x=A^{-1}\Delta b,
\]

se tiene

\[
    \|\Delta x\|
    \leq
    \|A^{-1}\|\,\|\Delta b\|.
\]

Además,

\[
    b=Ax
    \quad\Longrightarrow\quad
    \|b\|\leq\|A\|\|x\|.
\]

Combinando ambas desigualdades,

\[
    \frac{\|\Delta x\|}{\|x\|}
    \leq
    \|A\|\|A^{-1}\|
    \frac{\|\Delta b\|}{\|b\|}
    =
    \cond(A)
    \frac{\|\Delta b\|}{\|b\|}.
\]

Para la cota inferior, de

\[
    \Delta b=A\Delta x
\]

resulta

\[
    \|\Delta b\|
    \leq
    \|A\|\|\Delta x\|,
\]

mientras que

\[
    x=A^{-1}b
    \quad\Longrightarrow\quad
    \|x\|\leq\|A^{-1}\|\|b\|.
\]

Por lo tanto,

\[
    \frac{\|\Delta x\|}{\|x\|}
    \geq
    \frac{1}{\cond(A)}
    \frac{\|\Delta b\|}{\|b\|}.
\]
\end{proof}

La interpretación fundamental es que $\cond(A)$ actúa como un factor
máximo de amplificación del error relativo en los datos.


\subsubsection*{Perturbación de la matriz}

Consideremos ahora

\[
    Ax=b,
    \qquad
    (A+E)\widetilde{x}=b,
\]

y definamos

\[
    \Delta x=\widetilde{x}-x.
\]

Restando ambas ecuaciones,

\[
    A\Delta x+E\widetilde{x}=0,
\]

por lo que

\[
    \Delta x=-A^{-1}E\widetilde{x}.
\]

En consecuencia,

\[
    \boxed{
    \frac{\|\Delta x\|}
         {\|\widetilde{x}\|}
    \leq
    \|A^{-1}\|\,\|E\|
    =
    \cond(A)
    \frac{\|E\|}{\|A\|}.
    }
\]

Definamos

\[
    \delta
    =
    \cond(A)\frac{\|E\|}{\|A\|}.
\]

\begin{teorema}[Cota relativa para perturbaciones pequeñas]
Si

\[
    \delta<1,
\]

entonces $A+E$ es invertible y

\[
    \boxed{
    \frac{\|\Delta x\|}{\|x\|}
    \leq
    \frac{\delta}{1-\delta}.
    }
\]
\end{teorema}

\begin{proof}
De la desigualdad anterior,

\[
    \|\Delta x\|
    \leq
    \delta\|\widetilde{x}\|
    =
    \delta\|x+\Delta x\|.
\]

Por desigualdad triangular,

\[
    \|\Delta x\|
    \leq
    \delta\|x\|
    +
    \delta\|\Delta x\|.
\]

Luego,

\[
    (1-\delta)\|\Delta x\|
    \leq
    \delta\|x\|,
\]

y como $\delta<1$,

\[
    \frac{\|\Delta x\|}{\|x\|}
    \leq
    \frac{\delta}{1-\delta}.
\]
\end{proof}


\begin{observacion}
El condicionamiento es una propiedad del problema matemático
$Ax=b$, no del algoritmo utilizado para resolverlo.

Una matriz puede estar mal condicionada aunque se utilice un algoritmo
numéricamente estable. La estabilidad del algoritmo constituye un
concepto diferente y será estudiada por separado.
\end{observacion}


\subsection*{43. Distancia al Conjunto de Matrices Singulares}

Sea

\[
    \mathcal{S}
    =
    \{B\in\R^{n\times n}: B\text{ es singular}\}.
\]

El número de condición también posee una interpretación geométrica:
mide qué tan cerca se encuentra una matriz invertible de perder la
invertibilidad.

\begin{teorema}[Distancia a la matriz singular más cercana]
Sea $A\in\R^{n\times n}$ invertible y considérese una norma matricial
inducida. Entonces

\[
    \boxed{
    \inf_{B\in\mathcal{S}}
    \|A-B\|
    =
    \frac{1}{\|A^{-1}\|}.
    }
\]

Dividiendo por $\|A\|$,

\[
    \boxed{
    \frac{1}{\cond(A)}
    =
    \inf_{B\in\mathcal{S}}
    \frac{\|A-B\|}{\|A\|}.
    }
\]
\end{teorema}

Este resultado muestra que una matriz tiene número de condición grande
precisamente cuando se encuentra, en términos relativos, cerca del
conjunto de matrices singulares.

En particular, si para una familia $A_\varepsilon$ existe una familia
de matrices singulares $B_\varepsilon$ tal que

\[
    \frac{\|A_\varepsilon-B_\varepsilon\|}
         {\|A_\varepsilon\|}
    \longrightarrow0,
\]

entonces

\[
    \boxed{
    \cond(A_\varepsilon)\longrightarrow\infty.
    }
\]
\subsection*{Factorización LU}

La factorización $LU$ permite resolver un sistema lineal sin calcular explícitamente la inversa de la matriz. La eliminación de Gauss puede interpretarse como una sucesión de multiplicaciones por matrices triangulares inferiores elementales.

Para $i>j$, sea $m_{ij}$ el multiplicador utilizado para eliminar la entrada situada en la posición $(i,j)$. Una matriz elemental de eliminación tiene la forma
\[
    M_j
    =I-
    \begin{pmatrix}
    0&\cdots&0&0&\cdots&0\\
    \vdots&&\vdots&\vdots&&\vdots\\
    0&\cdots&m_{j+1,j}&0&\cdots&0\\
    0&\cdots&m_{j+2,j}&0&\cdots&0\\
    \vdots&&\vdots&\vdots&&\vdots
    \end{pmatrix},
\]
con entradas no nulas únicamente debajo de la diagonal en la columna $j$. Su inversa se obtiene cambiando el signo de dichos multiplicadores.

\begin{teorema}[Factorización LU]
Sea $A\in\R^{n\times n}$. Si la eliminación de Gauss puede realizarse sin intercambios de filas y sin utilizar pivotes nulos, entonces existen matrices $L$ y $U$ tales que
\[
    A=LU,
\]
donde $L$ es triangular inferior con diagonal unitaria y $U$ es triangular superior. La matriz $L$ contiene los multiplicadores empleados durante la eliminación.
\end{teorema}

Si $L_1,\ldots,L_{n-1}$ denotan las matrices triangulares inferiores elementales que representan las operaciones inversas de eliminación, entonces
\[
    U=L_{n-1}\cdots L_2L_1A,
\]
y, por consiguiente,
\[
    A=L_1^{-1}L_2^{-1}\cdots L_{n-1}^{-1}U.
\]
La matriz
\[
    L=L_1^{-1}L_2^{-1}\cdots L_{n-1}^{-1}
\]
es triangular inferior con unos en la diagonal, por lo que se obtiene $A=LU$.

\begin{observacion}
La condición de que no aparezcan pivotes nulos depende del orden de las filas. Cuando esto no se cumple, puede ser necesario intercambiar filas antes de continuar la eliminación.
\end{observacion}

\begin{teorema}[Factorización LU con pivoteo]
Sea $A\in\R^{n\times n}$ una matriz no singular. Entonces existe una matriz de permutación $P$ tal que
\[
    PA=LU,
\]
donde $L$ es triangular inferior con diagonal unitaria y $U$ es triangular superior. Las permutaciones de filas utilizadas durante la eliminación se incorporan en $P$.
\end{teorema}

Una matriz de permutación satisface
\[
    P^{-1}=P^{\mathsf T}.
\]
Si durante la eliminación se utilizan sucesivamente las permutaciones $P_1,\ldots,P_{n-1}$, la permutación total puede escribirse como el producto
\[
    P=P_{n-1}\cdots P_2P_1,
\]
con el orden de los factores determinado por la sucesión concreta de intercambios de filas.

\subsection*{Resolución del sistema mediante $LU$}

Una vez calculada la factorización, el sistema
\[
    Ax=b
\]
se transforma, en ausencia de pivoteo, en
\[
    LUx=b.
\]
Introduciendo $y=Ux$ o, de forma más habitual, $y$ como solución de la etapa triangular inferior, se resuelven sucesivamente
\[
    Ly=b,
    \qquad
    Ux=y.
\]
En el caso con pivoteo, $PA=LU$ y por tanto
\[
    LUx=Pb,
\]
de modo que primero se calcula $Pb$, luego se resuelve $Ly=Pb$ y finalmente $Ux=y$.

La ventaja fundamental es que resolver sistemas triangulares es mucho más sencillo y barato que calcular $A^{-1}$ de manera explícita. Además, una vez factorizada $A$, la misma factorización puede reutilizarse para varios términos independientes.

\subsection*{Factorización de Cholesky}

La factorización de Cholesky es una versión especialmente favorable de la factorización triangular para matrices simétricas definidas positivas.

Recordemos que una matriz simétrica $A\in\R^{n\times n}$ es \emph{definida positiva} si
\[
    x^{\mathsf T}Ax>0
    \qquad\text{para todo }x\neq0.
\]

\begin{teorema}[Teorema de Cholesky]
Una matriz $A\in\R^{n\times n}$ admite una factorización
\[
    A=LL^{\mathsf T},
\]
donde $L$ es triangular inferior con diagonal positiva, si y solo si $A$ es simétrica y definida positiva.
\end{teorema}

\begin{proof}
Si $A=LL^{\mathsf T}$, entonces
\[
    A^{\mathsf T}=(LL^{\mathsf T})^{\mathsf T}=LL^{\mathsf T}=A,
\]
de modo que $A$ es simétrica. Además, si $x\neq0$,
\[
    x^{\mathsf T}Ax
    =x^{\mathsf T}LL^{\mathsf T}x
    =(L^{\mathsf T}x)^{\mathsf T}(L^{\mathsf T}x)
    =\norm{L^{\mathsf T}x}_2^2>0,
\]
ya que $L$ es invertible. Por lo tanto, $A$ es definida positiva.

El recíproco es el teorema de existencia de Cholesky: toda matriz simétrica definida positiva admite una única factorización de esa forma cuando se exige que la diagonal de $L$ sea positiva.
\end{proof}

\begin{proposicion}[Fórmulas de Cholesky]
Sea $A=(a_{ij})$ simétrica definida positiva y sea $L=(\ell_{ij})$ triangular inferior. Los elementos de $L$ pueden calcularse sucesivamente mediante
\[
    \ell_{kk}
    =
    \sqrt{a_{kk}-\sum_{j=1}^{k-1}\ell_{kj}^2},
\]
y, para $i>k$,
\[
    \ell_{ik}
    =
    \frac{1}{\ell_{kk}}
    \left(
        a_{ik}-\sum_{j=1}^{k-1}\ell_{ij}\ell_{kj}
    \right).
\]
\end{proposicion}

Estas fórmulas se obtienen simplemente comparando las entradas de $A=LL^{\mathsf T}$. En particular, para $i\neq k$,
\[
    a_{ik}=\sum_{j=1}^{k}\ell_{ij}\ell_{kj},
\]
mientras que para la diagonal
\[
    a_{kk}=\sum_{j=1}^{k}\ell_{kj}^2.
\]

\subsection*{Resolución de un sistema por Cholesky}

Si $A=LL^{\mathsf T}$, el sistema
\[
    Ax=b
\]
puede escribirse como
\[
    LL^{\mathsf T}x=b.
\]
Definiendo $y$ mediante
\[
    Ly=b,
\]
queda después el sistema triangular superior
\[
    L^{\mathsf T}x=y.
\]
Así, la resolución se reduce a dos sustituciones triangulares, una hacia adelante y otra hacia atrás.

\begin{observacion}
La estructura de $L$ y $L^{\mathsf T}$ hace que no sea necesario calcular $A^{-1}$. En particular, para matrices simétricas definidas positivas, Cholesky aprovecha la simetría de la matriz y requiere trabajar esencialmente con la mitad triangular de sus elementos.
\end{observacion}

\subsection*{Relación entre las factorizaciones}

La factorización $LU$ y la factorización de Cholesky persiguen el mismo objetivo general: transformar un sistema lineal en sistemas triangulares más sencillos de resolver. La diferencia esencial es la estructura que se exige a la matriz original. La factorización $LU$ se aplica, con las condiciones de pivoteo correspondientes, a matrices generales; Cholesky requiere que la matriz sea simétrica definida positiva y produce una factorización especialmente estructurada,
\[
    A=LL^{\mathsf T}.
\]

\subsection*{Descomposición QR mediante Transformaciones de Householder}

\begin{teorema}
Dada una matriz $A \in \mathbb{C}^{m \times n}$, existen una matriz unitaria $Q$ y una matriz triangular superior $R$ tales que $A = QR$.
\end{teorema}

\begin{proof}
La demostración es de carácter constructivo y se fundamenta en la aplicación sucesiva de transformaciones unitarias, específicamente reflexiones de Householder. Comenzamos analizando la matriz por columnas, expresándola en bloques como $A = [A_1 | A_2 | \dots | A_n]$. Nuestro primer objetivo iterativo es hallar una matriz unitaria $U_1$ que, al premultiplicar a la primera columna $A_1$, anule todas sus componentes por debajo de la diagonal principal. Es decir, si tomamos el vector $x = A_1$, buscamos que la transformación resulte en $U_1 x = y$, donde el vector $y$ debe tener la forma $\beta e_1$.

Para evitar el fenómeno de cancelación catastrófica y asegurar algebraicamente que el producto interno entre los vectores involucrados sea un número real, el escalar $\beta$ se define de manera condicionada:
\begin{equation}
\beta = 
\begin{cases} 
-\frac{a_{11}}{|a_{11}|} \|A_1\|_2 & \text{si } a_{11} \neq 0 \\ 
-\|A_1\|_2 & \text{si } a_{11} = 0 
\end{cases}
\end{equation}
Con esta rigurosa elección, garantizamos en primer lugar la preservación de la norma, cumpliéndose que $\|x\|_2 = \|y\|_2$. Además, evaluando el producto interno entre $x$ e $y$, comprobamos que efectivamente pertenece a los reales:
\begin{equation}
\langle x, y \rangle = y^* x = (\beta e_1)^* x = \bar{\beta} a_{11} = \left( -\frac{\bar{a}_{11}}{|a_{11}|} \|A_1\|_2 \right) a_{11} = -|a_{11}| \|A_1\|_2 \in \mathbb{R}
\end{equation}

El procedimiento de construcción de esta primera matriz $U_1$ se sistematiza en una secuencia de cinco pasos. El paso 1 consiste en calcular el escalar $\beta$ tal como acabamos de definir. En el paso 2, inicializamos nuestros vectores de trabajo $x = A_1$ e $y = \beta e_1$. Para el paso 3, debemos definir el vector normal $V$ que generará el plano de reflexión: si $x \neq y$, se define como $V = \frac{x-y}{\|x-y\|_2}$; si resultaran iguales, simplemente $V = 0$. Avanzando al paso 4, ensamblamos la matriz de reflexión $U = I - V V^*$ (absorbiendo las constantes de proyección correspondientes en la definición de $V$). Por último, el paso 5 representa la aplicación directa de esta matriz $U$ sobre toda la matriz original $A$.

Antes de analizar la matriz resultante, es imperativo certificar la validez estructural de la herramienta que estamos utilizando.

\begin{proposicion}
La matriz extendida $U_i = \begin{pmatrix} I & 0 \\ 0 & I - V V^* \end{pmatrix}$, armada con la identidad de la dimensión adecuada en el bloque superior, es unitaria y hermítica.
\end{proposicion}
La hermiticidad, es decir $U_i^* = U_i$, se deduce de forma directa operando sobre la transpuesta conjugada: $(I - V V^*)^* = I^* - (V V^*)^* = I - V V^*$. Dado que la matriz representa geométricamente una reflexión plana, aplicarla dos veces consecutivas nos devuelve el estado original, por lo que $U_i^2 = I$. Esta propiedad involutiva, combinada con la hermiticidad, demuestra inmediatamente que la matriz es unitaria.

Asegurada la validez de la matriz, introducimos un resultado que formaliza el propósito geométrico de nuestra construcción.
\begin{lema}[Lema 2]
Para vectores distintos $x \neq y$, la transformación matricial cumple estrictamente que $(I - V V^*) x = y$.
\end{lema}
Al desarrollar analíticamente el miembro izquierdo mediante la fórmula de proyección ortogonal, obtenemos $(I - V V^*)x = x - \frac{(x-y)(x-y)^*}{\|x-y\|^2} (2x)$. Gracias a que demostramos previamente que el producto interno asociado es estrictamente real, el álgebra de los términos cruzados se cancela de manera perfecta, colapsando el resultado a $y$ y confirmando que la reflexión mapea el vector original al eje deseado.

Retomando el hilo principal en el paso 5, al multiplicar la matriz $A$ original por nuestra matriz unitaria $U_1$, la primera columna queda conformada únicamente por el elemento $\beta$ y ceros, dividiendo la estructura resultante en bloques:
\begin{equation}
U_1 A = \begin{pmatrix} \beta & z^* \\ 0 & W_1 \end{pmatrix}
\end{equation}

Para dar continuidad a la triangularización, aislamos la submatriz inferior definiendo $B_2 = W_1$. Sobre la primera columna de este nuevo bloque aplicamos exactamente los mismos pasos 1 a 5 descritos anteriormente, lo que nos proveerá de un nuevo vector director $v_2$. Con este vector se construye la matriz $U_2$, la cual debe incorporar un bloque identidad para no alterar la fila y columna que ya hemos procesado:
\begin{equation}
U_2 = \begin{pmatrix} 1 & 0 \\ 0 & I - v_2 v_2^* \end{pmatrix}
\end{equation}
Al aplicar $U_2$ sobre nuestra matriz en progreso obtenemos el siguiente paso de reducción:
\begin{equation}
U_2 U_1 A = \begin{pmatrix} 1 & 0 \\ 0 & I - v_2 v_2^* \end{pmatrix} \begin{pmatrix} \beta & z^* \\ 0 & W_1 \end{pmatrix} = \begin{pmatrix} \beta & z^* \\ 0 & (I - v_2 v_2^*) W_1 \end{pmatrix} = \begin{pmatrix} \beta_1 & * & * \\ 0 & \beta_2 & * \\ 0 & 0 & W_2 \end{pmatrix}
\end{equation}

Este mecanismo algorítmico es generalizable. Para un paso arbitrario $k$, habiendo definido una submatriz residual $W_k$, obtenemos su correspondiente vector de reflexión $V_{k+1}$ tras someter su primera columna a los pasos 1 al 5. La matriz unitaria del paso se define entonces como:
\begin{equation}
U_{k+1} = \begin{pmatrix} I_k & 0 \\ 0 & I_{m-(k-1)} - V_{k+1} V_{k+1}^* \end{pmatrix}
\end{equation}
donde el bloque $I_k$ representa la matriz identidad de dimensión $k$, cuya función es blindar las $k$ filas y columnas superiores ya estabilizadas. El bloque activo inferior opera exclusivamente sobre $W_k$, produciendo:
\begin{equation}
(I - V_{k+1} V_{k+1}^*) W_k = \begin{pmatrix} \beta_{k+1} & z_{k+1}^* \\ 0 & W_{k+1} \end{pmatrix}
\end{equation}

Visualizando el producto encadenado de todas estas transformaciones en su forma escalonada, tenemos:
\begin{equation}
U_{k+1} U_k \dots U_1 A = \begin{pmatrix} I_k & 0 \\ 0 & V_{k+1} \end{pmatrix} \begin{pmatrix} B_k & Z \\ 0 & W_k \end{pmatrix} = \begin{pmatrix} B_{k+1} & Z' \\ 0 & W_{k+1} \end{pmatrix}
\end{equation}

Este barrido iterativo por columnas culmina invariablemente en el paso $t = \min\{m, n-1\}$, momento en el que el algoritmo ya no encuentra más elementos subdiagonales que anular. El producto acumulado de todas las transformaciones unitarias aplicadas a la matriz original $A$ genera una nueva matriz $R$, la cual tiene, por construcción, una estructura estrictamente triangular superior:
\begin{equation}
U_t \dots U_1 A = R
\end{equation}
Como probamos en la Proposición, cada matriz de reflexión $U_i$ es unitaria y hermítica, por lo que coincide con su propia inversa. Despejando analíticamente la matriz $A$, trasladamos las transformaciones al miembro derecho:
\begin{equation}
A = (U_t \dots U_1)^{-1} R = U_1^* \dots U_t^* R
\end{equation}
Agrupando las matrices unitarias como $Q = U_1^* \dots U_t^*$, y apoyándonos en el teorema que dictamina que el producto de matrices unitarias da como resultado otra matriz unitaria, llegamos al enunciado final de la demostración: $A = QR$.
\end{proof}


\subsection*{44. Métodos Iterativos: Jacobi y Gauss--Seidel}

Consideremos el sistema lineal

\[
    Ax=b,
    \qquad
    A\in\mathbb{R}^{n\times n},
\]

donde $A$ es invertible y sus elementos diagonales son no nulos.

La idea de los métodos iterativos estacionarios consiste en construir
una sucesión de aproximaciones

\[
    x^{(0)},x^{(1)},x^{(2)},\ldots
\]

que, bajo ciertas condiciones, converja a la solución exacta $x$ del sistema.

En general, estos métodos pueden escribirse en la forma

\[
    \boxed{
    x^{(k+1)}=Bx^{(k)}+c
    }
\]

donde $B$ se denomina \emph{matriz de iteración}.

\subsubsection*{Descomposición de la matriz}

Escribimos

\[
    \boxed{
    A=D+L+U
    }
\]

donde

\[
D=
\begin{pmatrix}
a_{11} & 0      & \cdots & 0\\
0      & a_{22} & \cdots & 0\\
\vdots & \vdots & \ddots & \vdots\\
0      & 0      & \cdots & a_{nn}
\end{pmatrix},
\]

$L$ es la parte estrictamente triangular inferior de $A$, y $U$ es
la parte estrictamente triangular superior.

Por tanto,

\[
    (D+L+U)x=b.
\]


% ============================================================
\subsubsection*{Método de Jacobi}
% ============================================================

En el método de Jacobi se despeja la parte diagonal:

\[
    Dx=b-(L+U)x.
\]

La iteración queda definida por

\[
    Dx^{(k+1)}
    =
    b-(L+U)x^{(k)}.
\]

Como $D$ es invertible,

\[
    x^{(k+1)}
    =
    -D^{-1}(L+U)x^{(k)}
    +
    D^{-1}b.
\]

Por lo tanto,

\[
    \boxed{
    x^{(k+1)}
    =
    B_Jx^{(k)}+c_J
    }
\]

con

\[
    \boxed{
    B_J=-D^{-1}(L+U)
    }
\]

y

\[
    \boxed{
    c_J=D^{-1}b.
    }
\]


\paragraph{Forma componente a componente.}

El sistema

\[
    \sum_{j=1}^{n}a_{ij}x_j=b_i
\]

puede escribirse como

\[
    a_{ii}x_i
    =
    b_i-
    \sum_{\substack{j=1\\j\neq i}}^{n}a_{ij}x_j.
\]

Así, el método de Jacobi queda dado por

\[
    \boxed{
    x_i^{(k+1)}
    =
    \frac{1}{a_{ii}}
    \left(
    b_i-
    \sum_{\substack{j=1\\j\neq i}}^{n}
    a_{ij}x_j^{(k)}
    \right),
    \qquad i=1,\ldots,n.
    }
\]

La característica esencial del método es que todas las componentes
de $x^{(k+1)}$ se calculan utilizando únicamente componentes de
$x^{(k)}$.


% ============================================================
\subsubsection*{Método de Gauss--Seidel}
% ============================================================

En Gauss--Seidel se agrupan la diagonal y la parte triangular inferior:

\[
    (D+L)x=b-Ux.
\]

La iteración se define mediante

\[
    (D+L)x^{(k+1)}
    =
    b-Ux^{(k)}.
\]

Por tanto,

\[
    x^{(k+1)}
    =
    -(D+L)^{-1}Ux^{(k)}
    +
    (D+L)^{-1}b.
\]

Así,

\[
    \boxed{
    x^{(k+1)}
    =
    B_{GS}x^{(k)}+c_{GS}
    }
\]

donde

\[
    \boxed{
    B_{GS}=-(D+L)^{-1}U
    }
\]

y

\[
    \boxed{
    c_{GS}=(D+L)^{-1}b.
    }
\]


\paragraph{Forma componente a componente.}

La expresión de Gauss--Seidel es

\[
    \boxed{
    x_i^{(k+1)}
    =
    \frac{1}{a_{ii}}
    \left(
    b_i
    -
    \sum_{j<i}a_{ij}x_j^{(k+1)}
    -
    \sum_{j>i}a_{ij}x_j^{(k)}
    \right).
    }
\]

A diferencia de Jacobi, Gauss--Seidel utiliza inmediatamente las
componentes recién calculadas de la iteración $k+1$.


% ============================================================
\subsubsection*{Convergencia de un método iterativo estacionario}
% ============================================================

Sea

\[
    x^{(k+1)}=Bx^{(k)}+c,
\]

y sea $x$ un punto fijo del método, es decir,

\[
    x=Bx+c.
\]

Definiendo el error

\[
    e^{(k)}=x^{(k)}-x,
\]

se obtiene

\[
\begin{aligned}
    e^{(k+1)}
    &=
    x^{(k+1)}-x\\
    &=
    Bx^{(k)}+c-(Bx+c)\\
    &=
    B(x^{(k)}-x).
\end{aligned}
\]

Por tanto,

\[
    \boxed{
    e^{(k+1)}=Be^{(k)}
    }
\]

y, por recurrencia,

\[
    \boxed{
    e^{(k)}=B^ke^{(0)}.
    }
\]


\begin{teorema}[Criterio Fundamental de Convergencia]
El método iterativo

\[
    x^{(k+1)}=Bx^{(k)}+c
\]

converge para toda elección de $x^{(0)}$ si y sólo si

\[
    \boxed{
    \rho(B)<1.
    }
\]
\end{teorema}

Equivalentemente,

\[
    \rho(B)<1
    \quad\Longleftrightarrow\quad
    B^k\longrightarrow0.
\]

Por tanto,

\[
    \boxed{
    \rho(B_J)<1
    }
\]

es la condición necesaria y suficiente para la convergencia de Jacobi,
mientras que

\[
    \boxed{
    \rho(B_{GS})<1
    }
\]

es la condición necesaria y suficiente para la convergencia de
Gauss--Seidel.


% ------------------------------------------------------------
\subsubsection*{Criterio suficiente mediante normas}
% ------------------------------------------------------------

Como para cualquier norma matricial inducida se cumple

\[
    \rho(B)\leq\|B\|,
\]

se obtiene inmediatamente el siguiente criterio.

\begin{teorema}[Criterio mediante norma]
Si existe una norma matricial inducida tal que

\[
    \boxed{
    \|B\|<1,
    }
\]

entonces el método iterativo converge para toda condición inicial.
\end{teorema}

Esta condición es suficiente, pero no necesaria.


% ------------------------------------------------------------
\subsubsection*{Cota del error}
% ------------------------------------------------------------

Supongamos que

\[
    \|B\|\leq q<1.
\]

Entonces

\[
    \|e^{(k)}\|
    \leq
    q^k\|e^{(0)}\|.
\]

Además, puede obtenerse una estimación del error sin conocer la
solución exacta:

\[
    \boxed{
    \|x^{(k)}-x\|
    \leq
    \frac{q}{1-q}
    \|x^{(k)}-x^{(k-1)}\|.
    }
\]

Esta desigualdad proporciona un criterio práctico para detener la
iteración.


% ============================================================
\subsubsection*{Dominancia diagonal}
% ============================================================

\begin{definicion}[Dominancia Diagonal Estricta]
Una matriz $A=(a_{ij})$ es \emph{estrictamente diagonal dominante por
filas} si

\[
    \boxed{
    |a_{ii}|
    >
    \sum_{\substack{j=1\\j\neq i}}^{n}
    |a_{ij}|,
    \qquad
    i=1,\ldots,n.
    }
\]
\end{definicion}

\begin{teorema}[Convergencia por Dominancia Diagonal]
Si $A$ es estrictamente diagonal dominante por filas, entonces los
métodos de Jacobi y Gauss--Seidel convergen para toda elección de
$x^{(0)}$.
\end{teorema}

En particular,

\[
    A
    \text{ estrictamente diagonal dominante}
    \quad\Longrightarrow\quad
    \rho(B_J)<1
\]

y

\[
    A
    \text{ estrictamente diagonal dominante}
    \quad\Longrightarrow\quad
    \rho(B_{GS})<1.
\]

La dominancia diagonal constituye una condición suficiente, pero no
necesaria.


\subsubsection*{Gauss--Seidel para matrices simétricas definidas positivas}

\begin{teorema}
Sea $A\in\mathbb{R}^{n\times n}$ simétrica definida positiva, es decir,

\[
    A^T=A
\]

y

\[
    x^TAx>0
    \qquad
    \forall x\neq0.
\]

Entonces el método de Gauss--Seidel converge para toda elección de
$x^{(0)}$.
\end{teorema}

Por tanto,

\[
    \boxed{
    A\text{ simétrica definida positiva}
    \Longrightarrow
    \rho(B_{GS})<1.
    }
\]


% ============================================================
\subsubsection*{Comparación entre Jacobi y Gauss--Seidel}
% ============================================================

Cuando ambos métodos convergen, la velocidad asintótica de convergencia
está gobernada por el radio espectral de sus respectivas matrices de
iteración.

En términos generales, un radio espectral menor implica una
convergencia asintótica más rápida:

\[
    0\leq\rho(B_1)<\rho(B_2)<1
\]

indica que el método asociado a $B_1$ reduce el error más rápidamente
en el régimen asintótico.

Sin embargo, en general no existe una desigualdad universal entre

\[
    \rho(B_J)
    \qquad\text{y}\qquad
    \rho(B_{GS}).
\]


\begin{teorema}[Relación para matrices tridiagonales]
Sea $A$ una matriz tridiagonal con elementos diagonales no nulos y
considérese la partición natural

\[
    A=D+L+U.
\]

Bajo las hipótesis usuales de la ordenación natural, si $\mu\neq0$ es
un autovalor de la matriz de Jacobi $B_J$, entonces

\[
    \lambda=\mu^2
\]

es un autovalor correspondiente de la matriz de Gauss--Seidel
$B_{GS}$. En particular,

\[
    \boxed{
    \rho(B_{GS})
    =
    \rho(B_J)^2.
    }
\]
\end{teorema}

Por consiguiente, si

\[
    0<\rho(B_J)<1,
\]

entonces

\[
    \rho(B_{GS})
    =
    \rho(B_J)^2
    <
    \rho(B_J),
\]

por lo que Gauss--Seidel posee, en este caso, una convergencia
asintótica más rápida.


% ============================================================
\subsection*{45. Descomposición de Schur}
% ============================================================

La descomposición de Schur proporciona una forma triangular asociada a
cualquier matriz cuadrada. Su importancia radica en que permite estudiar
los autovalores de una matriz mediante una transformación que preserva
la geometría euclídea.


\begin{teorema}[Descomposición de Schur Compleja]
Sea

\[
    A\in\mathbb{C}^{n\times n}.
\]

Entonces existe una matriz unitaria $Q$ y una matriz triangular superior
$T$ tales que

\[
    \boxed{
    A=QTQ^*.
    }
\]

Equivalentemente,

\[
    \boxed{
    Q^*AQ=T.
    }
\]

Los elementos de la diagonal de $T$ son exactamente los autovalores de
$A$, contados con multiplicidad algebraica.
\end{teorema}


\paragraph{Interpretación.}

La matriz $Q$ representa un cambio de base ortonormal. Por tanto,
$A$ y $T$ son unitariamente semejantes y poseen los mismos autovalores:

\[
    \sigma(A)=\sigma(T).
\]

Como $T$ es triangular,

\[
    \boxed{
    \lambda_i(A)=t_{ii},
    \qquad
    i=1,\ldots,n.
    }
\]

En particular,

\[
    \boxed{
    \rho(A)
    =
    \max_{1\leq i\leq n}|t_{ii}|.
    }
\]


% ------------------------------------------------------------
\subsubsection*{Demostración de la Descomposición de Schur}
% ------------------------------------------------------------

La demostración puede realizarse por inducción sobre $n$.

Para $n=1$ el resultado es inmediato.

Supongamos ahora que el teorema es válido para matrices de tamaño
$(n-1)\times(n-1)$.

Como $A\in\mathbb{C}^{n\times n}$, existe al menos un autovalor
$\lambda_1$ y un autovector asociado $q_1$ tal que

\[
    Aq_1=\lambda_1q_1.
\]

Normalizamos $q_1$ de modo que

\[
    \|q_1\|_2=1
\]

y lo completamos hasta obtener una base ortonormal

\[
    \{q_1,q_2,\ldots,q_n\}.
\]

Definimos la matriz unitaria

\[
    Q_1=
    \begin{pmatrix}
    q_1&q_2&\cdots&q_n
    \end{pmatrix}.
\]

Entonces

\[
    Q_1^*AQ_1
    =
    \begin{pmatrix}
    \lambda_1 & *\\
    0 & A_1
    \end{pmatrix},
\]

donde

\[
    A_1\in\mathbb{C}^{(n-1)\times(n-1)}.
\]

Por hipótesis inductiva, existe una matriz unitaria $Q_2$ tal que

\[
    Q_2^*A_1Q_2=T_1,
\]

con $T_1$ triangular superior.

Tomando

\[
    \widetilde Q
    =
    \begin{pmatrix}
    1&0\\
    0&Q_2
    \end{pmatrix},
\]

se obtiene

\[
    \widetilde Q^*
    Q_1^*AQ_1
    \widetilde Q
    =
    \begin{pmatrix}
    \lambda_1 & *\\
    0&T_1
    \end{pmatrix},
\]

que es triangular superior.

Por tanto, tomando

\[
    Q=Q_1\widetilde Q,
\]

resulta

\[
    Q^*AQ=T,
\]

con $Q$ unitaria y $T$ triangular superior.
\qed


% ============================================================
\subsubsection*{Descomposición de Schur Real}
% ============================================================

Si

\[
    A\in\mathbb{R}^{n\times n},
\]

no siempre es posible triangularizar $A$ mediante una matriz ortogonal
real, pues $A$ puede poseer autovalores complejos.

Sin embargo, existe una matriz ortogonal $Q$ tal que

\[
    \boxed{
    A=QTQ^T,
    }
\]

donde $T$ es \emph{cuasi-triangular superior}: su diagonal está formada
por bloques de tamaño $1\times1$ y $2\times2$.

Los bloques $1\times1$ corresponden a autovalores reales, mientras que
los bloques $2\times2$ representan pares conjugados de autovalores
complejos.


% ============================================================
\subsubsection*{Relación con el Teorema Espectral}
% ============================================================

La descomposición de Schur existe para toda matriz cuadrada compleja,
sin necesidad de suponer simetría o normalidad.

Si además $A$ es normal,

\[
    A^*A=AA^*,
\]

la matriz triangular $T$ de la descomposición de Schur debe ser
diagonal.

Por tanto,

\[
    \boxed{
    A=Q\Lambda Q^*
    }
\]

con $\Lambda$ diagonal.

En consecuencia, el Teorema Espectral puede interpretarse como un caso
particular de la descomposición de Schur para matrices normales.

En el caso real simétrico,

\[
    A=A^T,
\]

se obtiene

\[
    \boxed{
    A=Q\Lambda Q^T,
    }
\]

donde $Q$ es ortogonal y $\Lambda$ es diagonal real.

\subsection*{Apéndice A: Herramientas Analíticas para Acotación de Errores}

En el Cálculo Numérico, evaluar las cotas del error local (mediante $\max|y''(t)|$ o derivadas superiores) y del amplificador geométrico (Constante de Lipschitz $L$) no requiere hallar el supremo exacto, sino construir un ``techo matemático inquebrantable'' válido en el dominio de simulación. Las siguientes herramientas analíticas permiten garantizar estos límites mitigando el costo de derivación y búsqueda de raíces.

\subsection*{A.1. Existencia y Ubicación (Weierstrass y Monotonía)}
\begin{itemize}
    \item \textbf{Teorema de Weierstrass:} Si una función $g(t)$ es continua en un intervalo cerrado y acotado $[a, b]$, entonces $g(t)$ alcanza obligatoriamente un valor máximo y un mínimo absoluto dentro de dicho intervalo. Esto garantiza teóricamente la existencia de la cota.
    \item \textbf{Criterio de Monotonía (El atajo analítico):} Si se demuestra mediante el signo de la derivada que $g'(t) \ge 0$ $\forall t \in [a, b]$, la función es estrictamente creciente. Por lo tanto, el máximo absoluto exigido por Weierstrass ocurre irremediablemente en el límite superior del intervalo, evaluándose directamente como $\max|g(t)| = |g(b)|$. (Análogamente, si es decreciente, ocurre en $a$).
\end{itemize}

\subsection*{A.2. Desigualdad Triangular (La Salida Pragmática)}
Para acotar sumas de funciones complejas sin calcular sus puntos críticos exactos, se emplea la desigualdad triangular:
\begin{equation}
    |A(t) + B(t)| \le |A(t)| + |B(t)|
\end{equation}
\textbf{El \textit{Trade-off}:} Genera un techo matemático superior al pico real de la curva, devolviendo una garantía de error más pesimista pero salvando el inmenso costo analítico de derivar e igualar a cero.

\subsection*{A.3. Cotas Trigonométricas y Exponenciales Clásicas}
\begin{itemize}
    \item \textbf{Cota de Amplitud Trigonométrica Exacta:} Debido al desfase de las ondas, el máximo absoluto de cualquier combinación lineal de la forma $A\sin(t) + B\cos(t)$ no es la suma de sus valores absolutos, sino la magnitud de su vector asociado:
    \begin{equation}
        \max |A\sin(t) + B\cos(t)| = \sqrt{A^2 + B^2}
    \end{equation}
    \item \textbf{Cota Exponencial de Decaimiento:} Para dominios donde $x \ge 0$, la función exponencial negativa está naturalmente acotada por la unidad: $e^{-x} \le 1$.
    \item \textbf{Identidad Exponencial de Funciones Positivas:} Si se demuestra que la trayectoria analítica $y(t) \ge 0$ durante la simulación, se garantiza físicamente que $e^{-y(t)} \le 1$. Además, la función $x e^{-x}$ alcanza un máximo absoluto de $1/e$ para $x \ge 0$.
\end{itemize}

\end{document}
